% =============================================================================
% DOCUMENTO FINAL — Proyecto de Transferencia de Calor
% Herramienta computacional educativa para el diseño de intercambiadores de
% calor de tubos y coraza — Método de Kern y norma TEMA Clase C
% =============================================================================
\documentclass[11pt,a4paper]{article}
\usepackage[margin=2.3cm, headsep=14pt]{geometry}
\usepackage[T1]{fontenc}
\usepackage[utf8]{inputenc}
\usepackage{xcolor}
\usepackage{booktabs}
\usepackage{longtable}
\usepackage{array}
\usepackage{fancyhdr}
\usepackage{titlesec}
\usepackage{caption}
\usepackage{enumitem}
\usepackage{microtype}
\usepackage{amsmath}
\usepackage{amssymb}
\usepackage{graphicx}
\usepackage{float}

\renewcommand{\tablename}{Tabla}
\renewcommand{\figurename}{Figura}
\renewcommand{\contentsname}{Contenido}

\definecolor{morado}{HTML}{5E2B8C}
\definecolor{acento}{HTML}{7E57C2}
\definecolor{seccion}{HTML}{EDE7F6}
\definecolor{gristexto}{HTML}{3B3648}

\titleformat{\section}{\Large\bfseries\color{morado}}{\thesection}{0.6em}{}[\vspace{2pt}{\color{acento}\titlerule}]
\titleformat{\subsection}{\large\bfseries\color{acento}}{\thesubsection}{0.5em}{}
\titleformat{\subsubsection}{\normalsize\bfseries\color{gristexto}}{\thesubsubsection}{0.5em}{}
\captionsetup{font=small, labelfont={bf,color=morado}, textfont={color=gristexto}}
\setlength{\parindent}{0pt}
\setlength{\parskip}{5pt}
\pagestyle{fancy}
\fancyhf{}
\fancyhead[L]{\small\color{morado}\textbf{Boh\'orquez $\cdot$ M\'arquez $\cdot$ Monta\~no $\cdot$ Mor\'e}}
\fancyhead[R]{\small\color{gristexto}Intercambiador de Tubos y Coraza $\cdot$ Kern / TEMA}
\fancyfoot[C]{\small\color{gristexto}\thepage}
\renewcommand{\headrulewidth}{0.4pt}

\emergencystretch=3em
\begin{document}

% -----------------------------------------------------------------------------
\begin{titlepage}
\centering
\vspace*{1.6cm}
{\color{acento}\rule{\linewidth}{2pt}}\\[1.0cm]
{\Huge\bfseries\color{morado} Herramienta computacional educativa\\[0.4em]
 para el dise\~no de intercambiadores de calor\\[0.4em]
 de tubos y coraza}\\[0.8cm]
{\Large\color{acento} M\'etodo de Kern (t\'ermico e hidr\'aulico)\\[0.2em]
 y norma TEMA Clase C (mec\'anico)}\\[0.5cm]
{\large\color{gristexto} Documento final del proyecto de Transferencia de Calor}\\[0.7cm]
{\large\color{gristexto}
\textbf{Boh\'orquez Bar\'on, Diego Nicol\'as} \quad \texttt{dibohorquezb@unal.edu.co}\\[0.25em]
\textbf{M\'arquez Ariza, Sara Camila} \quad \texttt{samarqueza@unal.edu.co}\\[0.25em]
\textbf{Monta\~no Ria\~no, Samuel} \quad \texttt{samontanor@unal.edu.co}\\[0.25em]
\textbf{Mor\'e D\'iaz, Carlos Andr\'es} \quad \texttt{cmore@unal.edu.co}}\\[1.0cm]
{\color{acento}\rule{\linewidth}{1pt}}\\[1.2cm]
{ \small
\begin{tabular}{rl}
\textbf{Software documentado:} & \texttt{IntercambiadorDeTubosYCoraza} (Python 3 / CustomTkinter)\\
\textbf{Funciones cubiertas:} & Sizing, rating, dise\~no mec\'anico y planos acotados\\
\textbf{Cobertura t\'ermica:} & 10 fluidos, 5 materiales, propiedades $\rho, c_p, k, \mu = f(T)$\\
\textbf{Cobertura TEMA:} & Tipos E/F/G/H/J/K/X y cabezales A/B/C/D/N/L/M/P/S/T/W\\
\textbf{Verificaci\'on:} & Suite \texttt{verificar.py}: 187 comprobaciones superadas\\
\textbf{Unidades:} & Sistema Internacional en todo el documento\\
\end{tabular}}
\vfill
{\small\color{gristexto} Todos los valores num\'ericos presentados provienen de la
 ejecuci\'on del c\'odigo fuente del proyecto sobre los once proyectos de ejemplo
 incluidos en \texttt{ejemplos/}, o de los documentos t\'ecnicos del propio
 proyecto. Ning\'un dato f\'isico ha sido estimado o supuesto por el redactor.}
\end{titlepage}

\tableofcontents
\newpage

% =============================================================================
\section{Resumen}
% =============================================================================

Este documento presenta el dise\~no, la formulaci\'on y la validaci\'on de una
herramienta computacional educativa para el c\'alculo de intercambiadores de
calor de tubos y coraza. El n\'ucleo t\'ermico implementa el m\'etodo can\'onico
de Kern: balance de energ\'ia, diferencia media logar\'itmica de temperatura
(LMTD) con factor de correcci\'on $F$ de Bowman, correlaciones de convecci\'on
para los lados tubos (Sieder--Tate y Hausen) y coraza (Kern), coeficiente
global referido al \'area exterior e iteraci\'on cerrada sobre $U$; el m\'odulo
de evaluaci\'on (\emph{rating}) resuelve las temperaturas de salida por
b\'usqueda de ra\'ices sobre el balance $q = U A F \,\Delta T_m$. El lado coraza se resuelve por
selecci\'on autom\'atica del m\'etodo: Bell--Delaware (factores $J_c$,
$J_l$, $J_b$ y $J_r$ sobre el banco ideal y ca\'ida de presi\'on en tres
componentes) mientras el Reynolds de coraza quede cubierto por la tabla
de Taborek ($1 \le Re_s \le 10^6$), y Kern fuera de ello; ambos
verificados contra la literatura. La selecci\'on es autom\'atica por
facilidad de uso, y cada proyecto registra el m\'etodo con el que se
calcul\'o. El
dimensionamiento mec\'anico sigue la norma TEMA Clase C (espejos, deflectores,
bridas, pernos, boquillas y cabezales). La interfaz gr\'afica, construida con
\texttt{customtkinter} bajo arquitectura Modelo--Vista--Controlador estricta, ofrece
sizing (F5), rating (F6), cinco vistas param\'etricas acotadas y persistencia
de proyectos en archivos \texttt{.bb}. Se reportan los resultados reales de los
once proyectos de ejemplo: ocho ejercicios resueltos de la literatura
(cuatro de sizing y cuatro de rating) y los tres del art\'iculo de Tan y Fok (2006).
Las comparaciones reproducen los Reynolds y las ca\'idas de presi\'on de la
literatura ($\Delta p_s$ de Janna $+0{,}1\%$; $\Delta p_t$ del
Ejemplo~7.4 de Kern $+0{,}1\%$; $h_o$ de Hewitt $-0{,}00\%$). La
consistencia bidireccional sizing--rating arroja errores relativos de carga de
$0{,}05\%$ a $0{,}44\%$. La suite de verificaci\'on autom\'atica
(187 comprobaciones) cierra sin fallas.

\textbf{Palabras clave:} intercambiador de calor de tubos y coraza, m\'etodo de
Kern, norma TEMA Clase C, dise\~no t\'ermico, software educativo, transferencia
de calor.

% =============================================================================
\section{Abstract}
% =============================================================================

This document presents the design, formulation and validation of an
educational computer tool for shell-and-tube heat exchanger calculation. The
thermal core implements Kern's classical method: energy balance, log-mean
temperature difference (LMTD) with Bowman's correction factor $F$, convection
correlations for the tube side (Sieder--Tate and Hausen) and the shell side
(Kern), the overall heat transfer coefficient referenced to the outer surface,
and a closed iteration on $U$; the rating module solves the outlet
temperatures by root finding on the balance
$q = U A F \,\Delta T_m$. Mechanical design follows the TEMA Class~C standard
(tubesheets, baffles, flanges, bolts, nozzles and heads). The graphical
interface, built with \texttt{customtkinter} under a strict Model--View--Controller
architecture, provides sizing (F5), rating (F6), five dimensioned parametric
views and project persistence in \texttt{.bb} files. Real results for the three included Tan and Fok (2006) projects ---and the
five service cases studied during development--- are reported: duties from $567{,}8$ to $5020{,}2$ kW,
overall coefficients from $548{,}2$ to $1312{,}3$ W/m\textsuperscript{2}K and
areas from $23{,}5$ to $141{,}5$ m\textsuperscript{2}. The two-way sizing--rating
consistency yields duty relative errors between $0{,}10\,\%$ and $0{,}70\,\%$
and $U$ errors below $0{,}01\,\%$. The automated verification suite
(138 checks) passes completely.

\textbf{Keywords:} shell-and-tube heat exchanger, Kern's method, TEMA Class~C
standard, thermal design, educational software, heat transfer.

% =============================================================================
\section{Introducci\'on}
% =============================================================================

El dise\~no de intercambiadores de calor de tubos y coraza es un contenido
central de la asignatura de Transferencia de Calor y, a la vez, uno de los
problemas m\'as ricos desde el punto de vista did\'actico: exige combinar
balances de energ\'ia, mecanismos de convecci\'on, resistencias t\'ermicas en
serie, criterios hidr\'aulicos y restricciones normativas de fabricaci\'on en un
\'unico procedimiento iterativo. La literatura especializada en ense\~nanza de
ingenier\'ia ha se\~nalado reiteradamente que las herramientas
computerizadas de apoyo al dise\~no favorecen esa integraci\'on; en particular,
Tan y Fok \cite{tanfok} describieron una herramienta computacional educativa
para el dise\~no de intercambiadores de calor de tubos y coraza que integra
el an\'alisis termohidr\'aulico con el dise\~no mec\'anico seg\'un la norma
TEMA y genera planos t\'ecnicos cotados. Su n\'ucleo t\'ermico sigue el
m\'etodo de Kern ---LMTD con el factor $F$ de los diagramas TEMA,
coeficientes de convecci\'on en el interior de los tubos y en el lado coraza
y ca\'idas de presi\'on en ambos lados---; el programa, escrito en Java con
orientaci\'on a objetos, se valid\'o frente a un problema de referencia del
texto de Incropera y deWitt, con errores m\'aximos del $2{,}2\,\%$ en sus
Tablas~1 (t\'ermica) e~2 (hidr\'aulica). El alcance de esa herramienta es
deliberadamente did\'actico: \textbf{no} optimiza de forma autom\'atica
---el usuario ajusta las entradas hasta alcanzar un dise\~no adecuado--- y
deja como trabajo futuro la estimaci\'on de costos, la lista de materiales y
las clases TEMA~B y~R. Ese desarrollo sirve de referencia conceptual para el
presente trabajo, que replica su planteamiento en Python con una interfaz
m\'as rica, verificaci\'on autom\'atica y los tres proyectos del propio
art\'iculo incluidos como ejemplos.

El objetivo del proyecto ha sido construir un software educativo
de nivel universitario, escrito en Python~3 con interfaz \texttt{customtkinter} (paleta de contraste morado, carmes\'i y naranja), que
calcule intercambiadores de calor de tubos y coraza en sus tres dimensiones:
\begin{enumerate}[leftmargin=1.4em]
\item \textbf{Dimensionamiento t\'ermico (\emph{sizing})} por el m\'etodo de
Kern: el usuario fija la geometr\'ia de los tubos y el programa itera sobre el
coeficiente global $U$ hasta obtener el \'area, el n\'umero de tubos y el
di\'ametro de coraza necesarios.
\item \textbf{Evaluaci\'on de un equipo existente (\emph{rating})}: con
$n_{tt}$ y $D_s$ conocidos, el programa itera sobre las temperaturas de salida
hasta predecir el calor realmente transferido.
\item \textbf{Dise\~no mec\'anico} seg\'un la norma TEMA Clase C y
\textbf{generaci\'on de planos} param\'etricos acotados.
\end{enumerate}

El alcance es deliberadamente acotado y expl\'icito: fluidos
\textbf{monof\'asicos} (l\'iquidos y gases), sin cambio de fase, dentro de las
ventanas de validez de temperatura declaradas para cada fluido; correlaciones
de Kern para el lado coraza y Sieder--Tate/Hausen para el lado tubos;
dimensionamiento mec\'anico seg\'un las tablas de la Secci\'on~5 del est\'andar
TEMA (Clase C). Las limitaciones del modelo se documentan en la
Secci\'on~\ref{sec:supuestos} y no deben interpretarse como defectos: cuando el
programa rechaza un dise\~no lo hace con un mensaje estructurado que explica el
principio f\'isico o la cl\'ausula normativa involucrada.

Este documento se organiza como sigue. La Secci\'on~\ref{sec:marco} desarrolla
el marco te\'orico (transferencia en coraza--tubos, m\'etodo de Kern y norma
TEMA). La Secci\'on~\ref{sec:desarrollo} presenta la metodolog\'ia, el sistema
completo de ecuaciones, los algoritmos de c\'alculo y el manual de usuario de la
interfaz. La Secci\'on~\ref{sec:verificacion} reporta la verificaci\'on con
ocho ejercicios de la literatura y la Secci\'on~\ref{sec:tanfok-analisis} el an\'alisis del caso de Tan y Fok. La
consistencia sizing--rating y el comportamiento t\'ermico e hidr\'aulico.
Finalmente, la Secci\'on~\ref{sec:conclusiones} recoge las conclusiones y el
Anexo~A contiene fragmentos representativos del c\'odigo fuente.

% =============================================================================
\section{Marco te\'orico}
\label{sec:marco}
% =============================================================================

\subsection{Transferencia de calor en intercambiadores de tubos y coraza}

Un intercambiador de tubos y coraza consta de un haz de tubos alojado dentro de
una coraza cil\'indrica. Una corriente circula por el interior de los tubos
(lado tubos) y la otra por el anillo entre el haz y la coraza (lado coraza), y
ambas intercambian calor a trav\'es de la pared de los tubos. Placas
deflectoras (\emph{baffles}) segmentadas obligan al fluido de la coraza a
atravesar el haz en flujo cruzado, lo que aumenta el coeficiente de convecci\'on
exterior y, al mismo tiempo, sostiene mec\'anicamente los tubos.

El transporte de calor en r\'egimen estacionario se describe como una
sucesi\'on de resistencias t\'ermicas en serie referidas al \'area exterior del
tubo:
convecci\'on interna, ensuciamiento interno, conducci\'on a trav\'es de la pared
cil\'indrica, ensuciamiento exterior y convecci\'on externa. Como las
corrientes no est\'an a temperatura uniforme, la fuerza motriz t\'ermica es una
diferencia media de temperaturas; para un arreglo en contracorriente esa media
es logar\'itmica y, para arreglos mixtos (varias pasadas por tubos o por
coraza), se corrige con el factor $F$ de Bowman. La potencia transferida se
escribe entonces
\begin{equation}
q = U_o \, A_o \, F \, \Delta T_m ,
\label{eq:q}
\end{equation}
donde $U_o$ es el coeficiente global sucio referido al \'area exterior $A_o$.

El flujo de la coraza es el fen\'omeno m\'as complejo del equipo: combina
flujo cruzado entre deflectores, fugas luz--coraza y tubo--deflector, y
corrientes de derivaci\'on. El m\'etodo de Kern lo representa mediante un
\'unico flujo cruzado ideal con un di\'ametro equivalente, lo que constituye
una simplificaci\'on controlada y adecuada para la ense\~nanza.

\subsection{El m\'etodo de Kern}

El procedimiento de Kern \cite{kern} resuelve dos problemas complementarios:

\begin{itemize}[leftmargin=1.4em]
\item \textbf{Dimensionamiento (sizing):} dados los caudales, las temperaturas
de entrada, la temperatura de salida de una corriente y la geometr\'ia de los
tubos ($d_o$, espesor, paso $P_T$, longitud $L$, n\'umero de pasos, arreglo),
se calcula el \'area $A$, el n\'umero de tubos $n_{tt}$ y el di\'ametro de
coraza $D_s$. El procedimiento es iterativo porque $U$ depende de las
velocidades, \'estas del \'area de flujo y el \'area de flujo del n\'umero de
tubos que a su vez depende de $U$.
\item \textbf{Evaluaci\'on (rating):} conocida la geometr\'ia real del equipo
($n_{tt}$, $D_s$), se predicen las temperaturas de salida y la potencia
transferida resolviendo la ecuaci\'on $q(T_{c,out}) - U A F \Delta T_m(T_{c,out}) = 0$.
\end{itemize}

El lado tubos cubre los tres reg\'imenes con correlaciones espec\'ificas:
soluci\'on de Graetz con entrada t\'ermica de Hausen (1959)\cite{hausen} en laminar y la
correlaci\'on de Gnielinski (1976)\cite{gnielinski} en transici\'on y turbulento; el lado coraza
usa la correlaci\'on de Kern turbulento y su rama laminar con dependencia de
Graetz ($D_e/B$). Las correlaciones que resuelven los reg\'imenes laminar y de
transici\'on ---Hausen y Gnielinski--- se verificaron contra la documentaci\'on
t\'ecnica consultada en l\'inea (gu\'ia de transferencia de calor de
neqsim/equinor\cite{guia}, secciones de flujo laminar, transici\'on y turbulento). El
programa se\~nala el r\'egimen resultante y reporta una banda de
incertidumbre estimada sobre $U$ (Secci\'on~\ref{sec:incertidumbre}). El
m\'etodo de Kern es un m\'etodo de dise\~no simplificado: no modela las fugas ni
la corriente de derivaci\'on de la coraza como lo hace el procedimiento de
Bell--Delaware, y su precisi\'on t\'ipica para fines de ingenier\'ia de proceso
es de un orden de magnitud del 10--20\,\% en el coeficiente de la coraza. Ese
es justamente su valor docente: un procedimiento completo, conservador y
trazable manualmente.

\subsection{Norma TEMA y Clase C}

La norma de la Tubular Exchanger Manufacturers Association \cite{tema}
estandariza la construcci\'on de estos equipos. La Figura N-1.2 del est\'andar
designa cada equipo con tres letras, \textbf{[cabezal frontal][coraza][cabezal
trasero]}:

\begin{center}
{\small
\begin{tabular}{@{}p{3.2cm}p{11.4cm}@{}}
\toprule
\textbf{Posici\'on} & \textbf{Tipos soportados por el software} \\
\midrule
Cabezal frontal & A (canal + cubierta), B (bonete), C (canal integral con
espejo, haz extra\'ible), D (cierre especial de alta presi\'on), N (canal
integral con espejo fijo) \\
Coraza & E (una pasada), F (dos pasadas con deflector longitudinal),
G (flujo partido), H (doble partido), J (dividido), K (caldera),
X (flujo cruzado) \\
Cabezal trasero & L/M/N (espejo fijo), P (flotante empaquetado),
S (flotante con anillo de retenci\'on), T (pull-through), U (haz en U),
W (espejo flotante sellado) \\
\bottomrule
\end{tabular}}
\end{center}

El est\'andar define tres clases de servicio ---R, C y B--- que fijan los
tolerancias, materiales y espesores admisibles. Este proyecto trabaja con la
\textbf{Clase C} (servicios comerciales de severidad moderada), cuyas tablas de
dimensionamiento mec\'anico (CB-3.1.3, RCB-4.3, CB-4.4.1, RCB-4.5.2, CB-4.7.1,
C-7.1.1, C-11.1 y RCB-4.6.1 de la Secci\'on~5) son las empleadas por el
m\'odulo de dise\~no mec\'anico. Las designaciones completas de los tipos
soportados y los diagramas correspondientes se reproducen en la vista de
cat\'alogo de la interfaz (Secci\'on~\ref{sec:manual}).

% =============================================================================
\section{Desarrollo: metodolog\'ia, ecuaciones y manual de usuario}
\label{sec:desarrollo}
% =============================================================================

\subsection{Arquitectura de c\'alculo}

El software separa estrictamente la f\'isica de la interfaz (arquitectura
Modelo--Vista--Controlador). El c\'odigo vive en cuatro m\'odulos:
\texttt{src/core/} (base de datos termof\'isica), \texttt{src/thermo/}
(\texttt{sizing.py} y \texttt{rating.py}), \texttt{src/mechanical/}
(\texttt{design.py}) y \texttt{src/gui/} (interfaz); \texttt{controller.py} es
el \'unico puente entre la vista y los m\'odulos de c\'alculo. La GUI no
contiene ni una l\'inea de f\'isica: incluso las ventanas de velocidad
recomendada que resalta en \'ambar se importan desde \texttt{sizing.py}, de modo
que el resaltado nunca puede contradecir al motor de c\'alculo.

El flujo de un c\'alculo completo es: (1) validaci\'on de entradas y de
ventanas de temperatura por punto; (2) balance de energ\'ia; (3) lazo iterativo
principal (sobre $U$ en sizing, sobre $T_{c,out}$ en rating); (4) cierre con
$n_{tt}$ entero y sobre-\'area; (5) dise\~no mec\'anico TEMA Clase C a partir de
los resultados t\'ermicos; (6) presentaci\'on y planos.

\subsection{Balance de energ\'ia}

En r\'egimen estacionario, sin p\'erdidas al ambiente ni generaci\'on interna y
con calores espec\'ificos evaluados a la temperatura media de cada corriente,
\begin{equation}
q = \dot m_h \, c_{p,h} \, (T_{h,in} - T_{h,out})
  = \dot m_c \, c_{p,c} \, (T_{c,out} - T_{c,in}).
\label{eq:balance}
\end{equation}
El programa exige que \textbf{exactamente una} temperatura de salida quede
vac\'ia: esa es la inc\'ognita de la ecuaci\'on~\eqref{eq:balance}. Las
temperaturas se cargan en $^\circ$C y los caudales en kg/s. Cada punto de
temperatura (entradas y salidas) se valida contra la ventana de validez del
fluido; por ejemplo, el etanol se rechaza a $95\,^\circ$C aunque su temperatura
media quede dentro de su rango $0$--$75\,^\circ$C, y el agua l\'iquida se limita
a $100\,^\circ$C (ebullici\'on a 1~atm).

\subsection{Propiedades termof\'isicas en funci\'on de la temperatura}

Las propiedades $\rho$, $c_p$, $k$ y $\mu$ se eval\'uan a la temperatura media
$T_m = (T_{in}+T_{out})/2$ mediante correlaciones emp\'iricas de alta fidelidad
(IAPWS para agua y vapor, DIPPR/Perry\'s para l\'iquidos org\'anicos,
correlaciones de referencia para gases). La viscosidad de pared $\mu_w$, que
aparece en la correcci\'on $\varphi = (\mu/\mu_w)^{0{,}14}$, se eval\'ua a la
temperatura de pared estimada,
\begin{equation}
T_w = \frac{h_{io}\,T_{t,m} + h_o\,T_{s,m}}{h_{io} + h_o},
\qquad h_{io} = h_i \, \frac{d_i}{d_o},
\end{equation}
con lazo interno acotado (hasta 40 pasos, convergencia en $10^{-9}$ sobre
$\varphi$). La Tabla~\ref{tab:fluidos} resume la cobertura de la base de datos.

\begin{table}[htbp]
\centering
\caption{Fluidos de la base termof\'isica: ventana de validez y resistencia de
ensuciamiento asignada.}
\label{tab:fluidos}
{\small
\begin{tabular}{@{}llll@{}}
\toprule
\textbf{Fluido} & \textbf{Ventana de validez ($^\circ$C)} &
\textbf{$R_f$ (m$^2$K/W)} & \textbf{Fuente de las correlaciones} \\
\midrule
Agua l\'iquida      & $0$ -- $100$   & $1{,}761 \times 10^{-4}$ & IAPWS \\
Vapor de agua       & $100$ -- $600$ & $8{,}81 \times 10^{-5}$  & IAPWS \\
Aire                & $-73$ -- $727$ & $3{,}522 \times 10^{-4}$ & Gas ideal + Sutherland \\
Etanol              & $0$ -- $75$    & $1{,}761 \times 10^{-4}$ & Perry / DIPPR \\
Metanol             & $0$ -- $60$    & $1{,}761 \times 10^{-4}$ & Perry / DIPPR \\
Etilenglicol        & $0$ -- $197$   & $3{,}522 \times 10^{-4}$ & REFPROP / DIPPR \\
Benceno             & $10$ -- $80$   & $1{,}761 \times 10^{-4}$ & NIST \\
Tolueno             & $-20$ -- $100$ & $1{,}761 \times 10^{-4}$ & NIST \\
n-Hexano            & $-20$ -- $60$  & $1{,}761 \times 10^{-4}$ & NIST \\
Amon\'iaco (l\'iquido) & $-40$ -- $50$ & $1{,}761 \times 10^{-4}$ & ASHRAE (saturado) \\
\bottomrule
\end{tabular}}
\end{table}

Los cinco materiales de construcci\'on disponibles son cobre
($k_m = 386{,}0$ W/m$\cdot$K, $\varepsilon = 1{,}5\ \mu$m), acero al carbono
($60{,}5$ W/m$\cdot$K, $\varepsilon = 46\ \mu$m), acero inoxidable AISI 304
($14{,}9$ W/m$\cdot$K), lat\'on Admiralty ($111{,}0$ W/m$\cdot$K) y aluminio
($237{,}0$ W/m$\cdot$K). La validaci\'on num\'erica de la base de datos frente
Sus conductividades y densidades provienen de tablas est\'andar a unos
$300$ K, y la rugosidad distingue tubo estirado ($\varepsilon = 1{,}5\ \mu$m)
de acero comercial ($\varepsilon = 46\ \mu$m). En resumen, la procedencia de
la base termof\'isica es: agua l\'iquida por IAPWS-IF97\cite{iapws}; vapor de agua y aire
por las tablas A.4 de Incropera\cite{incropera} (con correlaciones de Sutherland para $\mu$ y
$k$ del aire); alcoholes, hidrocarburos y glicol por el
\emph{Perry's Chemical Engineers' Handbook} (octava edici\'on)\cite{perry} y
correlaciones tipo DIPPR; y amon\'iaco l\'iquido por ASHRAE\cite{ashrae}/Perry's
(l\'iquido saturado). Los factores de ensuciamiento siguen la tabla TEMA
RGP-T-2.4. La validaci\'on num\'erica de la base de datos frente a valores
tabulados de referencia (NIST, CoolProp/REFPROP) reporta
desviaciones m\'aximas del $1{,}05\,\%$ en $\rho$ y $c_p$, $4{,}36\,\%$ en $k$
y $3{,}10\,\%$ en $\mu$ \cite{informe}.

\subsection{LMTD y factor de correcci\'on $F$}

Para contracorriente ideal,
\begin{equation}
\Delta T_m = \frac{\Delta T_1 - \Delta T_2}{\ln(\Delta T_1/\Delta T_2)},
\qquad
\Delta T_1 = T_{h,in} - T_{c,out},\quad
\Delta T_2 = T_{h,out} - T_{c,in}.
\end{equation}
Cuando $\Delta T_1 = \Delta T_2$ (cociente de capacidades $R = 1$) la
expresi\'on es la indeterminaci\'on $0/0$; el l\'imite por regla de L'H\^opital
es exactamente $\Delta T_1$, y el programa lo implementa con tolerancia
relativa $|\Delta T_1 - \Delta T_2| \le 10^{-6}\max(\Delta T_1,\Delta T_2)$,
devolviendo $0{,}5(\Delta T_1+\Delta T_2)$. El factor $F$ de Bowman corrige la
desviaci\'on respecto del contracorriente puro; se define con
\begin{equation}
P = \frac{T_{c,out}-T_{c,in}}{T_{h,in}-T_{c,in}}, \qquad
R = \frac{T_{h,in}-T_{h,out}}{T_{c,out}-T_{c,in}},
\end{equation}
y para el arreglo 1--2 (una pasada de coraza, $n$ pasadas de tubos) adopta la
forma cl\'asica
\begin{equation}
F = \frac{\sqrt{R^2+1}}{R-1}\,
\frac{\ln\!\big[(1-P)/(1-PR)\big]}
{\ln\!\left[\dfrac{2/P - 1 - R + \sqrt{R^2+1}}
{2/P - 1 - R - \sqrt{R^2+1}}\right]} ,
\end{equation}
con la rama $R \to 1$ tratada anal\'iticamente en el c\'odigo. El programa
aplica $F = 1$ (contracorriente puro) para una pasada de tubos y para los
tipos de coraza G, H, J, K y X; para el tipo F (dos pasadas de coraza en
serie) reduce $P$ por pasada antes de aplicar la expresi\'on anterior. Un
resultado $F < 0{,}75$ se rechaza como termodin\'amicamente inadmisible para el
arreglo elegido, y un $F$ inexistente (logaritmo no definido) se reporta como
\emph{cruce de temperaturas} con la acci\'on correctiva correspondiente.

\subsection{Coeficientes de pel\'icula $h_i$ y $h_o$}

\textbf{Lado tubos.} Con $\varphi = (\mu/\mu_w)^{0{,}14}$
(correcci\'on de Sieder--Tate (1936)\cite{sieder}) y $Gz = (d_i/L)\, Re_t\, Pr$
(n\'umero de Graetz), el n\'umero de Nusselt interno se calcula por ramas
seg\'un $Re_t = G_t d_i / \mu$:
\begin{align}
Re_t < 2100 \; &: \quad
Nu = 3{,}66 + \frac{0{,}0668\, Gz}{1 + 0{,}04\, Gz^{2/3}}\varphi
\quad \text{(Hausen, laminar)},\\[2pt]
2400 \le Re_t < 5\cdot 10^6 \; &: \quad
Nu = \frac{(f/8)(Re_t - 1000) Pr}{1 + 12{,}7\sqrt{f/8}\,(Pr^{2/3}-1)}
\Big[1 + (d_i/L)^{2/3}\Big] \varphi \quad \text{(Gnielinski)},\\[2pt]
Re_t \ge 5\cdot 10^6 \; &: \quad Nu = 0{,}023 \, Re_t^{0{,}8} Pr^{n}
\qquad \text{(Dittus--Boelter)}, \qquad n = 0{,}4/0{,}3.
\end{align}
La forma de Dittus--Boelter tiene dos exponentes seg\'un la direcci\'on del flujo de calor (Incropera, ec.~6.61):
\[
n = \begin{cases}
0{,}4 & \text{si el fluido se calienta } (T_w > T_m),\\[2pt]
0{,}3 & \text{si el fluido se enfr\'ia } (T_w < T_m),
\end{cases}
\]
La direcci\'on ya est\'a contenida en $n$, por lo que esta forma
no incorpora $\varphi$: aplicar ambas correcciones duplicar\'ia el efecto. En la franja $2100 \le Re_t < 2400$ las ramas de Hausen y Gnielinski se
mezclan en escala logar\'itmica, de modo que $Nu(Re)$ queda continua sin
alterar las f\'ormulas cl\'asicas fuera de la franja y sin saltos que
desestabilicen el lazo iterativo. El factor de fricci\'on de Gnielinski es el
de Petukhov, $f = (0{,}790\, \ln Re_t - 1{,}64)^{-2}$. La forma de Hausen es
la soluci\'on de Graetz con entrada t\'ermica: recupera $Nu \to 3{,}66$
(pared isot\'ermica, tubo largo) cuando $Gz \to 0$ y el asintoto de entrada
$Nu \approx 1{,}67\, Gz^{1/3}$ para $Gz$ grande, sin el piso duro que ten\'ia
la forma cl\'asica $1{,}86\,Gz^{1/3}$. El coeficiente interno es
$h_i = Nu \, k / d_i$

\textbf{Lado coraza.} Para deflectores segmentados al $25\,\%$,
\begin{align}
Re_s \ge 2000 \; &: \quad Nu_s = 0{,}36 \, Re_s^{0{,}55} Pr^{1/3} \varphi
\qquad \text{(Kern turbulento)},\\[2pt]
Re_s < 100 \; &: \quad Nu_s = \max\!\big(1{,}86 \, (Re_s\, Pr\, D_e/B)^{1/3}
\varphi,\; 3{,}66\big) \quad \text{(laminar, Sieder--Tate)},
\end{align}
y entre $100$ y $2000$ (transici\'on) se mezcla en escala log--log entre la
rama laminar y la turbulenta, ambas evaluadas en el mismo $Re_s$, de modo que
$Nu_s$ es continua en $100$ y en $2000$ y conserva la dependencia $D_e/B$ en
toda la banda. El coeficiente
exterior es $h_o = Nu_s \, k / D_e$. En la rama laminar interviene el espaciado
de deflectores $B$ a trav\'es de $D_e/B$, porque el desarrollo t\'ermico depende
de la longitud del canal de flujo cruzado.

\subsection{Geometr\'ia de flujo, di\'ametro equivalente y conteo de tubos}

El \'area de flujo cruzado en el eje de la coraza (Kern, ec.~7.1) y el \'area de
flujo por paso de tubos son
\begin{equation}
A_s = \frac{D_s \, (P_T - d_o) \, B}{P_T} \cdot \frac{1}{n_{corazas}},
\qquad
a_t = \frac{n_{tt} \, \pi \, d_i^2}{4 \, n_{pasos}},
\end{equation}
donde $n_{corazas} = 2$ para corazas TEMA F y $1$ para el resto. La
conservaci\'on de masa fija las velocidades: $G_t = \dot m_t / a_t$,
$v_t = G_t/\rho_t$, $G_s = \dot m_s / A_s$, $v_s = G_s/\rho_s$; nunca se imponen
velocidades que violen la continuidad. El di\'ametro equivalente de Kern se
basa en el perimetro mojado de transferencia:
\begin{equation}
D_e =
\begin{cases}
\dfrac{4\left(\dfrac{\sqrt{3}}{4}P_T^2 - \dfrac{\pi}{8}d_o^2\right)}
{\pi \, d_o / 2}, & \text{arreglo triangular (30°/60°)},\\[10pt]
\dfrac{4\left(P_T^2 - \dfrac{\pi}{4}d_o^2\right)}{\pi \, d_o}, &
\text{arreglo cuadrado (45°/90°)} .
\end{cases}
\end{equation}

El n\'umero de tubos se estima con la correlaci\'on de conteo de Serth
\cite{serth} con las constantes de arreglo $C_L$ ($0{,}87$ triangular, $1{,}00$
cuadrado) y de partici\'on de pasos $C_{TP}$ ($0{,}93$ para 1 paso, $0{,}90$
para 2, $0{,}85$ para 4 o m\'as):
\begin{equation}
D_s = \text{holgura} + P_T \sqrt{\frac{4}{\pi}\frac{C_L}{C_{TP}} n_{tt}} ,
\label{eq:ds}
\end{equation}
donde la holgura reserva el espacio haz--coraza que exige el cabezal trasero
elegido (de $0{,}011$ m en espejo fijo a $0{,}092$ m en cabezal T). La
correlaci\'on de conteo preliminar usa la forma
$D_s \approx \text{holgura} + 0{,}637\sqrt{C_L/C_{TP}}\sqrt{A\,(P_T/d_o)^2 d_o/L}$.

\subsection{Coeficiente global $U$}

Referido al \'area exterior y con resistencias de ensuciamiento $R_{f,i}$ y
$R_{f,o}$ (Kern, ec.~6.11):
\begin{equation}
\frac{1}{U_o} =
\frac{d_o}{d_i \, h_i}
+ R_{f,i}\,\frac{d_o}{d_i}
+ \frac{d_o \ln(d_o/d_i)}{2 \, k_m}
+ R_{f,o}
+ \frac{1}{h_o}.
\label{eq:u}
\end{equation}
La resistencia de ensuciamiento se asigna con precedencia estricta: valor
provisto por el usuario, valor TEMA del fluido o, en su defecto, la mitad de la
resistencia del material. $U$ resulta siempre del inverso de una suma de
resistencias positivas; un valor no f\'isico detiene el c\'alculo con un mensaje
explicativo.

\subsection{Ca\'idas de presi\'on}
\label{sec:dp}

Lado tubos (Kern, ecs.~7.43 y 7.45/7.46): la ca\'ida total es la suma de
fricci\'on y retornos en los cabezales,
\begin{equation}
\Delta p_t = n_{pasos} \, f_t \, \frac{L}{d_i}\,
\frac{\rho_t v_t^2}{2}\,\frac{1}{\varphi_t^{\,m}}
+ 4 \, n_{pasos} \, \frac{\rho_t v_t^2}{2},
\label{eq:dpt}
\end{equation}
con $m = 0{,}14$ en ambos r\'egimenes: es la correcci\'on de Kern por 
viscosidad de pared, aplicada sin introducir exponentes laminares sin 
respaldo en la fuente. El factor $f_t$ es $64/Re$ en laminar y la correlaci\'on
expl\'icita de Haaland (aproximaci\'on de Colebrook--White para tubos rugosos)
en turbulento.

Lado coraza (Kern, ec.~7.44):
\begin{equation}
\Delta p_s = \frac{n_{corazas} \, f_s \, G_s^2 \, D_s \, (N_b + 1)}
{2 \, \rho_s \, D_e \, \varphi_s},
\qquad
N_b = \left\lfloor \frac{L}{B} + 10^{-9} \right\rfloor - 1 ,
\label{eq:dps}
\end{equation}
con el factor de Kern $f_s = \exp(0{,}576 - 0{,}19 \ln Re_s)$ (empalmado a una
ley hiperb\'olica para $Re_s < 400$). El t\'ermino $10^{-9}$ dentro del piso
elimina errores de truncamiento en coma flotante (p. ej. $0{,}6/0{,}2 =
2{,}999\ldots$). En las boquillas se verifica el criterio de vibraci\'on de
TEMA RCB-4.6.1, $\rho V^2 \le 2232{,}25$ Pa ($1500$ lb$_m$/ft$\cdot$s$^2$), de
donde la velocidad m\'axima admisible $V_{max} = \sqrt{2232{,}25/\rho}$.


\subsection{Verificaci\'on de las ca\'idas de presi\'on frente a la
literatura}

Las ecuaciones anteriores reproducen literalmente las ecuaciones 7.44, 7.45 y
7.46 de Kern (1950). En particular, la p\'erdida por retorno de los tubos se
calcula con la convenci\'on de Kern: \emph{cuatro cabezas de velocidad por
paso} (cita literal del libro: \emph{``the return loss~\ldots accounted for
by allowing four velocity heads per pass''}). La literatura posterior discrepa
de esa cifra, y esa es la raz\'on por la que las ca\'idas de presi\'on de
este trabajo pueden resultar entre un $11,5$ y un $37,8\%$ por encima de
las que devuelven otros m\'etodos: Serth (2003) recomienda $(2 n_p - 1)$ cabezas
en total, Frank $2{,}5$ por paso, Bell--Mueller $3$ por paso, HEXTRAN $2$ por
paso y Saunders (1988) $K = 0{,}9$ (un paso) o $1{,}6$ (dos o m\'as) por paso.
Para el lado coraza existen dos ajustes publicados de la misma gr\'afica de
Kern: $f_s = \exp(0{,}576 - 0{,}19\, \ln Re_s)$ (este trabajo) y
$f_s = 1{,}728\, Re_s^{-0{,}188}$ (Lemos et al., Univ. de Oklahoma); la
diferencia entre ambos ajustes es de solo $0{,}6$--$1{,}7\%$.

Como contraste num\'erico con la literatura publicada, aplicando las
ecuaciones de este trabajo al equipo del art\'iculo de Tan y Fok con la
velocidad que ese art\'iculo impone ($v = 0{,}6$ m/s) y un tubo 
estirado ($\varepsilon = 1{,}5\,\mu$m) se obtiene 
$\Delta p_t = 1582$ Pa frente a sus $1975{,}59$ Pa publicados 
(desviaci\'on del $-19{,}9\%$). La diferencia no es de geometr\'ia sino 
del factor de fricci\'on: el valor publicado implica $f_t = 0{,}035$, que 
corresponde a una rugosidad $\varepsilon \approx 57\,\mu$m; con esa 
rugosidad el programa reproduce los 1975 Pa dentro del $0{,}1\%$. 
La ejecuci\'on sobre el mismo equipo con la velocidad real que impone 
su geometr\'ia ($0{,}14$ m/s) da $\Delta p_t = 109$ Pa: la discrepancia aparente con la tabla del
art\'iculo no es del m\'etodo sino de la velocidad de dise\~no impuesta,
ya discutida en la Secci\'on~\ref{sec:tanfok}.


\subsection{Criterio de la velocidad en tubos y el n\'umero de Reynolds}

Los ejercicios de aplicaci\'on del curso (Cuadro~1 y ec.~6 del documento de
ejercicios basado en Tan y Fok) definen el n\'umero de Reynolds del lado
tubos con la velocidad \emph{impuesta} $U_m = 0{,}59997$ m/s:
$Re_t = \rho_t U_m d_i / \mu_t$, que con las propiedades del enunciado da
$Re_t \approx 8{,}6\cdot 10^3$. Este trabajo, en cambio, sigue el criterio
de Kern: la velocidad se obtiene por \emph{continuidad} con el haz propuesto,
de modo que para el equipo del art\'iculo (257 tubos en un paso) la velocidad
real es $0{,}14$ m/s y el Reynolds resultante $Re_t \approx 1{,}8\cdot
10^3$. La diferencia entre ambos Reynolds no es un error de c\'alculo sino el
criterio de velocidad empleado: cada Reynolds es consistente con su velocidad.
El propio documento de ejercicios contempla las dos opciones ---su Cuadro~3
compara expresamente las filas \emph{``Dise\~no, $U_m$ por continuidad
($0{,}14$ m/s)''} y \emph{``Dise\~no, $U_m$ impuesta ($0{,}59997$
m/s)''}--- y su secci\'on 3.5 se\~nala que el m\'etodo del art\'iculo
\emph{impone} la velocidad sin ligarla al n\'umero de tubos. Al comparar
resultados con la literatura debe verificarse, por tanto, con qu\'e criterio
de velocidad fue calculado cada Reynolds.

\subsection{Algoritmo de dimensionamiento (sizing)}

El procedimiento implementado es el siguiente:
\begin{enumerate}[leftmargin=1.6em]
\item Balance de energ\'ia \eqref{eq:balance}: $q$ y la temperatura de salida
faltante.
\item LMTD y $F$ con las temperaturas completas.
\item \textbf{Lazo sobre $U$}: para cada $U$ de prueba se calcula
$A = q/(U F \Delta T_m)$, el n\'umero de tubos fraccionario $n_{tt,f} =
A/(\pi d_o L)$ y el $D_s$ tentativo \eqref{eq:ds}; con esa geometr\'ia se
obtienen $Re$, $Nu$, $h_i$, $h_o$ (con lazo de $\varphi$ de pared) y un
$U_{calc}$ nuevo.
\item La ecuaci\'on $U_{calc}(U) - U = 0$ se resuelve por \textbf{bisecci\'on
con bracket fijo} $[1, 5\times 10^3]$ W/m\textsuperscript{2}K (ampliado por
doblamiento si hiciera falta). Se eligi\'o bisecci\'on y no iteraci\'on directa
porque $U_{calc}(U) \sim C U^{\beta}$ con $\beta < 1$: la iteraci\'on directa
converge lento y su punto final depende de la semilla $U_{inicial}$, sobre todo
en laminar; con bracket fijo el resultado es el mismo con cualquier semilla.
Tolerancia relativa $10^{-4}$.
\item \textbf{Cierre con haz real}: $n_{tt} = \lceil A/(\pi d_o L)\rceil$
redondeado al m\'ultiplo superior de $n_{pasos}$ (reparto equitativo por paso);
se reeval\'ua todo el estado con el haz instalado, se recalcula
$A = q/(U_{real} F \Delta T_m)$ y se reporta la sobre-\'area
$(A_{real} - A)/A \times 100$.
\item Deflectores $N_b = \lfloor L/B \rfloor - 1$ y dise\~no mec\'anico.
\end{enumerate}

El espaciado de deflectores $B$ es dato de entrada, pero su ventana v\'alida
($D_s/5 \le B \le D_s$) depende de $D_s$, que es resultado: ese lazo se cierra
en dos pasadas manuales asistidas por el bot\'on
\textbf{$\frac{1}{2}\cdot D_s$} de la
interfaz, que coloca $B = D_s/2$ para recalcular (F5).

\subsection{Algoritmo de evaluaci\'on (rating) y par\'ametros $\varepsilon$--NTU}

Con la geometr\'ia fija, el rating resuelve la ra\'iz de
\begin{equation}
g(T_{c,out}) = q(T_{c,out}) - U \, A_o \, F \, \Delta T_m(T_{c,out}) = 0
\end{equation}
por un m\'etodo \textbf{secante/bisecci\'on} sobre un bracket $[T_{c,in},
T_{h,in}]$: se prueba el paso de la secante y, si cae fuera del $95\,\%$ central
del bracket, se usa el punto medio, conservando as\'i la garant\'ia de la
bisecci\'on con la velocidad de la secante. El criterio de parada es un bracket
m\'as angosto que $0{,}01\,^\circ$C (m\'aximo 60 iteraciones). Adem\'as del punto
de operaci\'on, se reportan los par\'ametros de contraste
\begin{equation}
\varepsilon = \frac{q}{q_{max}}, \qquad
NTU = \frac{U A_o}{C_{min}}, \qquad
C_r = \frac{C_{min}}{C_{max}},
\end{equation}
\'utiles para contrastar con soluciones $\varepsilon$--NTU.

\subsection{Dise\~no mec\'anico TEMA Clase C}

A partir del resultado t\'ermico ($D_s$, $n_{tt}$, $L$, $B$, presiones de
dise\~no, corrosi\'on y materiales) el m\'odulo mec\'anico dimensiona:
\textbf{coraza} (espesor $T_s$ seg\'un tablas TEMA con corrosi\'on, $D_o =
D_s + 2 T_s$); \textbf{espejo de tubos} ($L_{ts} = 0{,}1 \, D_s$, di\'ametros
l\'imite del haz $D_{otl}$ y $D_{ctl}$); \textbf{cubiertas de cabezal}
($T_c$, plana apernada para canal tipo A, toriesf\'erica para bonete);
\textbf{bridas} (ancho $f_w \approx 3\,(1{,}5\, d_B)$,
$D_F = D_o + 2 f_w + 0{,}004$ m, c\'irculo de pernos $D_{BC}$, n\'umero y paso
de pernos); \textbf{deflectores} ($N_B$, espesor tabulado, di\'ametro
$D_b$, altura de corte al $25\,\%$, luz sin soporte para verificaci\'on de
vibraci\'on); \textbf{varillas tensores} y \textbf{boquillas} (di\'ametro
m\'inimo por $\rho V^2 \le 2232{,}25$ Pa y selecci\'on de tuber\'ia comercial
NPS). El m\'odulo verifica tambi\'en el empaque por \'area sobre $D_{otl}$:
discrepancias menores al $5\,\%$ entre el conteo de tubos y el empaque se
reportan como advertencia, no como error.

\subsection{Unidades SI y supuestos del modelo}
\label{sec:supuestos}

\begin{enumerate}[leftmargin=1.6em]
\item Todas las ecuaciones se resuelven en \textbf{unidades SI}: m, kg, s, K
(o $^\circ$C en diferencias), W, Pa, W/m\textsuperscript{2}K. Las temperaturas
de entrada se cargan en $^\circ$C; las propiedades se eval\'uan en K internamente.
\item Flujo \textbf{monof\'asico} y permanente: no se modelan condensadores ni
evaporadores (la coraza K se calcula solo en su parte sensible).
\item Propiedades constantes a la temperatura media de cada coraza/corriente,
con $\mu_w$ para la correcci\'on $\varphi$.
\item Correlaciones de Kern: sin correcci\'on de fugas ni corriente de
derivaci\'on (Bell--Delaware); resultados orientativos en r\'egimen laminar y de
transici\'on (bandas de incertidumbre reportadas).
\item El balance de energ\'ia ignora p\'erdidas al ambiente y cambio de fase.
\item El agua l\'iquida se limita a $100\,^\circ$C y el amon\'iaco se modela como
l\'iquido saturado; no se cubren servicios presurizados de ebullici\'on.
\item El dise\~no mec\'anico sigue las tablas TEMA Clase C citadas; no sustituye
un c\'alculo de certificaci\'on ASME para servicio cr\'itico.
\end{enumerate}

\subsection{Manual de usuario}
\label{sec:manual}

La aplicaci\'on se lanza con \texttt{py src/gui/app.py} (o el atajo
\texttt{ejecutar.bat} en Windows) y no requiere dependencias externas ni
conexi\'on a internet. La ventana principal organiza el trabajo en cuatro
pesta\~nas: \textbf{Design}, \textbf{Analysis}, \textbf{Drawing} y
\textbf{Help}.

\subsubsection{Pesta\~na Design: formulario de entrada}

La interfaz est\'a construida con \texttt{customtkinter} y una paleta de
contraste \textbf{morado, carmes\'i y naranja}: controles redondeados,
tarjetas de selecci\'on, bloques de formulario blancos y botones de acci\'on
destacados. Lo primero que muestra la hoja de dise\~no es el
\textbf{selector de modo}
(paso~1): tres opciones destacadas ---\textbf{DIMENSIONAR un equipo nuevo}
(sizing), \textbf{EVALUAR un equipo existente} (rating) y \textbf{OTRAS}
(formulario completo, el men\'u cl\'asico)---, cada una con una ayuda que
enumera exactamente qu\'e datos necesita. Seg\'un el modo elegido, el
formulario muestra solo los bloques que corresponden, de modo que los datos
de sizing y de rating nunca se mezclan:

\begin{itemize}[leftmargin=1.4em]
\item \textbf{Bloques comunes:} corrientes de proceso (fluidos seleccionables
por men\'u desplegable --con la ventana de validez de temperatura visible
bajo cada uno--, temperaturas de entrada y de salida, caudales en kg/s y
asignaci\'on de cada fluido al lado tubos o coraza); geometr\'ia del tubo
($d_o$, espesor, paso $P_T$, longitud $L$, n\'umero de pasos, arreglo
triangular o cuadrado, material); coraza y deflectores (espaciado $B$ con el
bot\'on \textbf{$\frac{1}{2}\cdot D_s$}, corte de deflector); configuraci\'on
TEMA (tipo de coraza y cabezales, con la \textbf{designaci\'on TEMA}
--p.~ej.~\texttt{AEM}-- en vivo) y condiciones de dise\~no y materiales.
\item \textbf{Bloque A -- solo sizing:} la semilla de $U$ y los
\textbf{l\'imites permisibles de ca\'ida de presi\'on} ($\Delta p$ m\'axima
en tubos y en coraza) para la verificaci\'on hidr\'aulica; adem\'as se pide
dejar vac\'ia exactamente una «T salida» (la inc\'ognita del balance
$q = \dot m c_p \Delta T$).
\item \textbf{Bloque B -- solo rating:} la geometr\'ia del equipo existente
($n_{tt}$ y $D_s$); las dos «T salida» quedan vac\'ias porque las calcula el
programa.
\end{itemize}

El modo \textbf{OTRAS} muestra ambos bloques a la vez, con el formulario
completo tradicional. Los botones de la barra inferior ---o las teclas
\textbf{F5} (sizing) y \textbf{F6} (rating)--- lanzan el c\'alculo de forma
as\'incrona, sin congelar la ventana; si un proyecto de rating se calcula con
F5, el programa lo detecta y ejecuta la evaluaci\'on, indic\'andolo en la
barra de estado.

% ---- marcador de captura -----------------------------------------------------
\begin{figure}[H]
\centering
\includegraphics[width=0.92\linewidth]{Imagenes Nuevas/0_1_1_formulario_design.png}
\caption{Formulario de entrada (a) (pesta\~na Design).}
\label{fig:design}
\end{figure}

% ---- marcador de captura -----------------------------------------------------
\begin{figure}[H]
\centering
\includegraphics[width=0.92\linewidth]{Imagenes Nuevas/0_1_2_formulario_design.png}
\caption{Formulario de entrada (b) (pesta\~na Design).}
\label{fig:design}
\end{figure}

\paragraph{El bot\'on \texttt{½·$D_s$}.}
El espaciado de deflectores $B$ es un \emph{dato de entrada} del c\'alculo
t\'ermico, pero su ventana v\'alida ---$D_s/5 \le B \le D_s$
(TEMA RCB-4.5.1)--- depende de $D_s$, que es un \emph{resultado} del
dimensionamiento. Esa dependencia circular se cierra a mano en dos pasadas con
el bot\'on \texttt{½·$D_s$} situado junto al campo $B$:

\begin{enumerate}[leftmargin=1.6em]
\item calcule con \textbf{F5} usando un $B$ tentativo (p. ej. $250$ mm);
\item pulse \texttt{½·$D_s$}: el campo $B$ se rellena con la
mitad del $D_s$ del \'ultimo c\'alculo ($B = 0{,}5\,D_s$, el punto medio
de la recomendaci\'on $0{,}3\,D_s \le B \le 1\,D_s$) y la barra de
estado confirma el ajuste;
\item recalcule con \textbf{F5}. Dos pasadas bastan en la pr\'actica:
el $D_s$ converge y $B$ queda dentro de la ventana TEMA.
\end{enumerate}

Si el bot\'on se pulsa \emph{antes de} cualquier c\'alculo, no hay
$D_s$ que repartir y el programa lo explica en un di\'alogo, dejando el
campo intacto. El bot\'on no calcula nada por s\'i mismo: solo propone el
$B$; el c\'alculo sigue siendo siempre una acci\'on expl\'icita del
usuario (F5/F6). Su funcionamiento est\'a cubierto por la suite de
verificaci\'on autom\'atica.

\subsubsection{Pesta\~na Analysis: resultados}

Tras calcular, la pesta\~na Analysis presenta cuatro \'areas:
\begin{itemize}[leftmargin=1.4em]
\item \textbf{Thermal Design:} carga $q$, coeficiente global $U$, \'area
requerida e instalada, sobre-\'area, $n_{tt}$, $D_s$, $N_b$, LMTD, $F$,
temperaturas de salida, coeficientes $h_i$ y $h_o$, n\'umeros de Reynolds,
velocidades, di\'ametro equivalente, reg\'imenes de flujo y factores de
fricci\'on.

% ---- marcador de captura -----------------------------------------------------
\begin{figure}[H]
\centering
\includegraphics[width=0.92\linewidth]{Imagenes Nuevas/02_resultados_thermal.png}
\caption{Resultados t\'ermicos (pesta\~na Analysis).}
\label{fig:thermal}
\end{figure}

\item \textbf{Hydraulic Design:} ca\'idas de presi\'on lado tubos y coraza,
comparadas con los l\'imites permisibles cargados, y verificaci\'on de las
ventanas de velocidad recomendadas (tubos $0{,}9$--$2{,}4$ m/s; coraza
$0{,}3$--$1{,}5$ m/s) y del criterio $\rho V^2$ de las boquillas.

% ---- marcador de captura -----------------------------------------------------
\begin{figure}[H]
\centering
\includegraphics[width=0.92\linewidth]{Imagenes Nuevas/03_resultados_hydraulic.png}
\caption{Resultados hidr\'aulicos (pesta\~na Analysis).}
\label{fig:hydraulic}
\end{figure}

\item \textbf{Verificaci\'on del dise\~no:} un checklist de verificaciones con
estado \texttt{OK}/\texttt{AVISO} \textbf{adaptado al modo de trabajo}. En
\emph{sizing} se verifica la propuesta de equipo nuevo: sobre-\'area no
negativa y razonable, pitch ratio $P_T/d_o \ge 1{,}25$, espaciado de baffles
$0{,}2\,D_s \le B \le D_s$ (TEMA RCB-4.5.1), relaci\'on $L/D_s$ entre 3 y 15,
corte de deflector 15--45~\%, ca\'idas de presi\'on frente a los l\'imites
permisibles, ventanas de velocidad, $F \ge 0{,}75$ y validez de la
correlaci\'on de Kern. En \emph{rating} se verifica la evaluaci\'on del
equipo existente: cierre del balance $\lvert q - UAF\Delta T_m\rvert$,
efectividad $\varepsilon \le 95\,\%$ (distancia al pinch), segunda ley, y las
mismas verificaciones hidr\'aulicas compartidas.

% ---- marcador de captura -----------------------------------------------------
\begin{figure}[H]
\centering
\includegraphics[width=0.92\linewidth]{Imagenes Nuevas/03_1_resultados_desygn.png}
\caption{Verificación del diseño (pesta\~na Analysis).}
\label{fig:hydraulic}
\end{figure}

\item \textbf{Advertencias y notas educativas:} el r\'egimen detectado, la
banda de incertidumbre de $U$ si aplica, y los avisos de pitch ratio,
velocidades o empaque.
\end{itemize}

% ---- marcador de captura -----------------------------------------------------
\begin{figure}[H]
\centering
\includegraphics[width=0.92\linewidth]{Imagenes Nuevas/03_2_resultados_ADS.png}
\caption{Advertencias y notas (pesta\~na Analysis).}
\label{fig:hydraulic}
\end{figure}



\subsubsection{Pesta\~na Drawing: vistas param\'etricas}

La pesta\~na Drawing genera cinco vistas dibujadas sobre \texttt{Canvas} a
partir de las cotas del dise\~no mec\'anico (no son im\'agenes fijas):
\textbf{1. Corte longitudinal} del equipo con el s\'imbolo del tipo de coraza y
cabezales elegidos; \textbf{2. Distribuci\'on de tubos} (arreglo triangular o
cuadrado con $P_T$, $d_o$, $D_{otl}$ y $D_{ctl}$); \textbf{3. Placa deflectora}
(di\'ametro, altura de corte, espaciado y luz sin soporte); \textbf{4. Cabezal
y brida} (canal o bonete, espesor de cubierta, brida, c\'irculo de pernos y
boquillas); \textbf{5. Cat\'alogo TEMA (Fig. N-1.2)} con la selecci\'on actual
resaltada. Un clic sobre cualquier pieza muestra sus cotas en mil\'imetros en
el panel lateral.

% ---- marcadores de captura --------------------------------------------------
\begin{figure}[H]
\centering
\includegraphics[width=0.48\linewidth]{Imagenes Nuevas/04_corte_longitudinal.png}
\hspace{0.04\linewidth}
\includegraphics[width=0.48\linewidth]{Imagenes Nuevas/05_distribucion_tubos.png}
\caption{Vistas 1 (corte longitudinal) y 2 (distribuci\'on de tubos).}
\label{fig:vistas12}
\end{figure}

\begin{figure}[H]
\centering
\includegraphics[width=0.48\linewidth]{Imagenes Nuevas/06_baffle.png}
\hspace{0.04\linewidth}
\includegraphics[width=0.48\linewidth]{Imagenes Nuevas/07_cabezal_brida.png}
\caption{Vistas 3 (placa deflectora) y 4 (cabezal y brida).}
\label{fig:vistas34}
\end{figure}

\begin{figure}[H]
\centering
\includegraphics[width=0.92\linewidth]{Imagenes Nuevas/08_catalogo_tema.png}
\caption{Resultados hidr\'aulicos (pesta\~na Analysis).}
\label{fig:hydraulic}
\end{figure}

\subsubsection{Proyectos \texttt{.bb}: abrir, guardar y ejemplos}

El men\'u \textbf{File} ofrece Nuevo, Abrir, Guardar y Salir. Los proyectos se
guardan con extensi\'on \texttt{.bb} en formato JSON, con las claves
\texttt{formato} (valor \texttt{intercambiador-tubos-coraza}), \texttt{version}
(valor 1), \texttt{descripcion} y el bloque \texttt{inputs} con todos los
campos del formulario; el programa rechaza archivos con formato o versi\'on
incompatibles con un mensaje explicativo. La
opci\'on \textbf{File $\rightarrow$ Abrir ejemplo validado} carga cualquiera de
los tres proyectos verificados de la carpeta \texttt{ejemplos/}: los del
artículo de referencia de Tan y Fok (2006)---, que se
calculan solos.

\begin{figure}[H]
\centering
\includegraphics[width=0.75\linewidth]{Imagenes Nuevas/09_guardar_bb.png}
\caption{Di\'alogo de guardado de proyectos (\texttt{.bb}).}
\label{fig:guardar}
\end{figure}

\subsubsection{Mensajes de error educativos}

Cuando el modelo o la norma rechazan un dise\~no, la aplicaci\'on muestra una
ventana estructurada en cuatro preguntas ---\textbf{QU\'E PAS\'O} (el hecho, con
n\'umeros), \textbf{POR QU\'E} (el principio f\'isico o la cl\'ausula de la
norma), \textbf{C\'OMO LO SOLUCIONO} (la acci\'on concreta sobre el formulario)
y la \textbf{REFERENCIA}--- y sigue ejecut\'andose. Casos t\'ipicos: violaci\'on
de la segunda ley (el fluido caliente entra m\'as fr\'io que el fr\'io),
cruce de temperaturas imposible en contracorriente, fluido fuera de su ventana
de validez, agua por encima de $100\,^\circ$C, balance mal planteado (debe
quedar vac\'ia exactamente una temperatura de salida) y dise\~nos no
fabricables seg\'un TEMA (empaque, boquillas o luz sin soporte).

% ---- marcador de captura -----------------------------------------------------
\begin{figure}[H]
\centering
\includegraphics[width=0.80\linewidth]{Imagenes Nuevas/10_ejemplo_dialogo_error.png}
\caption{Ejemplo de mensaje de error educativo estructurado.}
\label{fig:error}
\end{figure}

% =============================================================================
\subsection{M\'etodo de Bell--Delaware (opci\'on de coraza)}
\label{sec:bell}

El programa permite calcular el lado coraza con el m\'etodo de Kern 
o con el m\'etodo de \textbf{Bell--Delaware} \cite{bell}. Ambos son \textbf{monof\'asicos}: Bell--Delaware 
no es un m\'etodo de cambio de fase, sino una descomposici\'on del flujo 
en cinco corrientes (cruzado ideal, fugas deflector--coraza y 
deflector--tubo, bypass haz--coraza y partici\'on de paso) que corrige el 
coeficiente del banco ideal:

\begin{equation}
h_o = h_{id}\,J_c\,J_l\,J_b\,J_r,\qquad
h_{id} = j_i\,c_p\,G_s\,Pr^{-2/3}\left(\frac{\mu}{\mu_w}\right)^{0{,}14},
\label{eq:hid}
\end{equation}

con $G_s = \dot m_s/S_m$ en el \'area de flujo cruzado, el 
Reynolds de coraza evaluado con el \textbf{di\'ametro exterior} ($Re_s = d_o G_s/\mu$; 
es una diferencia de definici\'on con Kern, que usa $D_e$) y el 
factor de Colburn ideal

\begin{equation}
j_i = a_1\left(\frac{1{,}33}{P_T/d_o}\right)^{a} Re_s^{a_2},\qquad
a = \frac{a_3}{1+0{,}14\,Re_s^{a_4}},
\label{eq:ji}
\end{equation}

cuyos coeficientes $a_1\ldots a_4$ (y $b_1\ldots b_4$ para el factor 
de fricci\'on $f_i$) dependen del rango de Reynolds y del \'angulo de 
arreglo y se toman de la tabla de Taborek \cite{taborek} recogida en la secci\'on 9.4.4 
de Kaka\c{c} y Liu \cite{kakac}. Los factores de correcci\'on son

\begin{align}
J_c &= 0{,}55 + 0{,}72 F_c, \qquad
J_l = 0{,}44(1-r_s) + \big[1-0{,}44(1-r_s)\big]e^{-2{,}2 r_{lm}},
\label{eq:jl}\\[2pt]
J_b &= e^{-C_{bh}F_{sbp}}, \qquad
J_r = 1 \ (Re_s>100),
\label{eq:jb}
\end{align}

y la ca\'ida de presi\'on se descompone en cruzado, ventana y 
extremos,

\begin{equation}
\Delta p_s = \Delta p_c + \Delta p_w + \Delta p_e,\qquad
\Delta p_c = \Delta p_{bi}(N_b-1)R_bR_l,
\label{eq:dpbd}
\end{equation}

con $\Delta p_{bi} = 2 f_i N_{cc} G_s^2/\rho$, 
y $\Delta p_e = \Delta p_{bi}(1+N_{cw}/N_{cc})R_bR_s$ 
($R_s = 2$ con espaciado uniforme). Las \'areas de fuga usan las 
holguras TEMA ($L_{bb} = 12{,}8+4{,}8D_s$ mm, $L_{sb} = 3{,}1+4D_s$ mm y 
$L_{tb} = 0{,}8$ mm para tubos de hasta 1 pulgada).

\textbf{Validaci\'on del m\'etodo en el programa.} Aplicado al equipo del 
Ejemplo~7.4 de Kern, los factores caen en los rangos publicados para 
intercambiadores bien dise\~nados: $J_c = 1{,}02$ (0,53--1,15), 
$J_l = 0{,}80$ (0,7--0,8), $J_b = 0{,}81$ (0,7--0,9), $J_r = 1{,}00$, con 
producto $J_cJ_lJ_bJ_r = 0{,}66$ (valor t\'ipico $\approx 0{,}6$); la 
correlaci\'on ideal resulta continua en los bordes de la tabla (salto 
m\'aximo 0,7\,\% en $j_i$ y 6,8\,\% en $f_i$). La comparaci\'on frente a Kern en los 
equipos de ejemplo se resume en la Tabla~\ref{tab:bd}.

\begin{table}[H]
\centering
\caption{Lado coraza con los dos m\'etodos: Kern frente a Bell--Delaware 
en los equipos de ejemplo.}
\label{tab:bd}
{\small
\begin{tabular}{@{}lrrrrrr@{}}
\toprule
\textbf{Caso} & $U_K$ & $U_{BD}$ & $h_{o,K}$ & $h_{o,BD}$ & $\Delta p_{s,K}$ & $\Delta p_{s,BD}$ \\
 & (W/m\textsuperscript{2}K) & (W/m\textsuperscript{2}K) & (W/m\textsuperscript{2}K) & (W/m\textsuperscript{2}K) & (Pa) & (Pa) \\
\midrule
R2 Kern 7.4 (rating) & 1430 & 1390 & 5707 & 5122 & 54\,224 & 24\,336 \\
R1 Janna (rating)    & 1126 & 1156 & 4299 & 4752 & 22\,504 & 16\,146 \\
R3 msubbu (rating)   & 783  & 742  & 1794 & 1594 & 29\,063 & 9\,671 \\
S4 Kern 7.4 (sizing) & 1454 & 1398 & 6171 & 5379 & 60\,787 & 34\,363 \\
S3 Janna (sizing)    & 1205 & 1246 & 4036 & 4436 & 20\,794 & 14\,275 \\
\bottomrule
\end{tabular}}
\end{table}

Las diferencias son las esperables del m\'etodo: $h_o$ de Bell--Delaware 
queda entre el 87 y el 111\,\% del de Kern (descuenta las fugas y el 
bypass que Kern agrupa) y $\Delta p_s$ resulta 30--60\,\% menor, porque 
elimina el car\'acter agrupado y conservador de Kern. En el caso de Tan y 
Fok, Bell--Delaware lleva el dise\~no a $U = 433$ W/m\textsuperscript{2}K 
y 272 tubos, frente a los 257 y el $U \approx 447$ que implica el 
art\'iculo: la desviaci\'on baja del 224\,\% de Kern al 
5,8\,\%, lo que confirma que la diferencia de ese caso estaba en 
el modelo de coraza y no en el programa. La banda de incertidumbre 
tambi\'en distingue el m\'etodo: con Bell--Delaware el aporte de la 
coraza baja de $\pm 15/20/25\%$ a $\pm 12/15/20\%$ seg\'un el 
r\'egimen (Secci\'on~\ref{sec:incertidumbre}).

\textbf{Selecci\'on autom\'atica del m\'etodo.} El m\'etodo 
del lado coraza se elige autom\'aticamente, por facilidad de uso: el 
programa aplica Bell--Delaware mientras el Reynolds de coraza quede 
dentro de la tabla de Taborek ($1 \le Re_s \le 10^6$) y vuelve a Kern 
fuera de ese rango, donde la extrapolaci\'on del gr\'afico de Kern es 
m\'as robusta. El m\'etodo realmente usado se registra en el resultado 
y dentro del proyecto \texttt{.bb}, de modo que cada tabla de este 
documento se reproduce abriendo el ejemplo correspondiente. Los once proyectos de ejemplo se conservaron con el m\'etodo de su 
propia fuente: los ocho ejercicios de la literatura con Kern, porque los 
valores publicados que verifican se calcularon con el m\'etodo de Kern, y 
compararlos con Bell--Delaware medir\'ia diferencias de m\'etodo, no 
errores del programa (en el Ejemplo~7.4 de Kern, Bell--Delaware da 
$U = 1390$ W/m\textsuperscript{2}K frente a los 1471 publicados, mientras 
que el m\'etodo de Kern reproduce 1430). Para dise\~no nuevo, en cambio, 
la selecci\'on autom\'atica aplica Bell--Delaware, que es el m\'etodo 
m\'as riguroso; la comparaci\'on de la Tabla~\ref{tab:bd} y el an\'alisis 
del caso de Tan y Fok (Secci\'on~7) cuantifican esa diferencia.

\section{Verificaci\'on del programa con ocho ejercicios de la literatura}
\label{sec:verificacion}
% =============================================================================

\subsection{M\'etodo y fuentes de la verificaci\'on}

Para demostrar que la herramienta calcula correctamente, se la someti\'o a
\textbf{ocho ejercicios resueltos y publicados} ---cuatro de dimensionamiento
(\emph{sizing}) y cuatro de evaluaci\'on (\emph{rating})--- tomados de la
literatura de dise\~no de intercambiadores. Se eligieron fuentes que publican
no solo cargas y temperaturas sino tambi\'en \textbf{n\'umeros de Reynolds
y ca\'idas de presi\'on}, de modo que la verificaci\'on alcanza el
c\'alculo hidr\'aulico completo y no solo el t\'ermico. Los ocho problemas
se cargaron como proyectos del programa (carpeta \texttt{ejemplos/}) y se
ejecutaron sin modificar el motor:

\begin{itemize}[leftmargin=1.4em]
\item \textbf{Janna} \cite{janna}, \emph{Design of Thermal Systems},
p.~251: agua destilada--agua cruda (rating con la geometr\'ia propuesta de
196 tubos y redimensionamiento del mismo servicio).
\item \textbf{Kern} \cite{kern}, Ejemplo~7.4: agua destilada--agua cruda
(rating del equipo disponible de 160 tubos y redimensionamiento del servicio).
\item \textbf{Sinnott y Towler} \cite{sinnott}
(Coulson \& Richardson, vol.~6), Ejemplo~12.1, y su verificaci\'on
publicada por \textbf{Subramanian} \cite{msubbu}: metanol--agua salobre
(dise\~no completo y verificaci\'on del equipo de 918 tubos).
\item \textbf{Kharaji} \cite{kharaji} (caso de Karimi et al.,
Tablas~4--6): metanol--agua de mar (dise\~no \'optimo PSO y
evaluaci\'on de su geometr\'ia).
\end{itemize}

\subsection{Ejercicios de dimensionamiento (sizing)}

La Tabla~\ref{tab:sizing8} compara el dise\~no publicado con el que propone
el programa a partir de la misma especificaci\'on de proceso (caudales,
temperaturas y geometr\'ia de tubos).

\begin{table}[htbp]
\centering
\caption{Dimensionamiento de cuatro servicios de la literatura: valores
publicados frente a los del programa.}
\label{tab:sizing8}
\footnotesize
\setlength{\tabcolsep}{4pt}
\begin{tabular}{@{}lrrrrrrrr@{}}
\toprule
 & \multicolumn{4}{c}{\textbf{Publicado}} &
\multicolumn{4}{c@{}}{\textbf{Programa}} \\
\cmidrule(lr){2-5}\cmidrule(l){6-9}
\textbf{Caso} & $q$ & $U$ & $n_{tt}$ & $D_s$ & $q$ & $U$ & $n_{tt}$ & $D_s$ \\
 & (kW) & (W/m\textsuperscript{2}K) & (--) & (m) & (kW) &
(W/m\textsuperscript{2}K) & (--) & (m) \\
\midrule
S1 Sinnott--T. 12.1 & 4340 & 738  & 918  & 0{,}894 & 4323 & 1147 & 408  & 0{,}571 \\
S2 Karimi (PSO)     & 4340 & 901  & 1238 & 0{,}811 & 4323 & 780  & 1442 & 0{,}804 \\
S3 Servicio Janna   & 1011 & 1277 & 196  & 0{,}438 & 1011 & 1205 & 268  & 0{,}486 \\
S4 Servicio Ej. 7.4 & 410  & 1471 & 160  & 0{,}400 & 410  & 1454 & 162  & 0{,}347 \\
\bottomrule
\end{tabular}
\end{table}

Las \textbf{cargas t\'ermicas coinciden en los cuatro casos dentro del $0{,}4\%$}, tal como corresponde a un balance de energ\'ia independiente del m\'etodo. Adem\'as, el sizing selecciona el n\'umero de pasos de tubos para que la velocidad quede dentro de la ventana recomendada (Kern la elige as\i; la longitud del tubo es un dato): con los pasos cargados en el formulario el punto de convergencia puede quedar con velocidades bajas, coeficientes bajos y por tanto demasiados tubos. El caso \textbf{S4} es el contraste m\'as exigente y tambi\'en el mejor: el programa propone $162$ tubos frente a los $160$ del equipo publicado ($+1{,}3\%$) con $U = 1454$ frente a $1471$ W/m\textsuperscript{2}K ($-1{,}1\%$). En \textbf{S3} el ajuste de pasos (de $2$ a $4$) lleva la velocidad a $1{,}31$ m/s y el dise\~no pide $268$ tubos con $U = 1205$ W/m\textsuperscript{2}K, un $5{,}6\%$ por debajo del $1277$ publicado: la diferencia restante frente a los $196$ tubos del enunciado es la convenci\'on del factor $F$ ($0{,}834$ en el arreglo 1-2) y la propia inconsistencia de esa fuente, cuyo $U\cdot A\cdot F\cdot LMTD$ no alcanza su propia carga. En \textbf{S1} el propio dise\~no publicado est\'a sobredimensionado (su $U$ calculado de $738$ triplica el supuesto de $600$ y su autor recomienda a\~nadir pasos).

\subsection{Ejercicios de evaluaci\'on (rating)}

La Tabla~\ref{tab:rating8} eval\'ua cuatro equipos publicados con sus
condiciones de operaci\'on. Es el contraste m\'as exigente: cada
Reynolds y cada ca\'ida de presi\'on se compara con el valor que la
fuente imprime.

\begin{table}[H]
\centering
\caption{Evaluaci\'on de cuatro equipos publicados: valores de la
literatura frente a los del programa (entre par\'entesis, la desviaci\'on).}
\label{tab:rating8}
\small
\begin{tabular}{@{}llrr@{}}
\toprule
\textbf{Magnitud} & \textbf{Fuente} & \textbf{Publicado} &
\textbf{Programa} \\
\midrule
\multicolumn{4}{@{}l@{}}{\itshape R1: Janna, agua destilada--agua cruda
(196 tubos)} \\
$T_{s,o}$ ($^\circ$C)     & Janna          & 32{,}11   & 33{,}67  \\
$T_{t,o}$ ($^\circ$C)     & Janna          & 31{,}11   & 29{,}26  \\
$U$ (W/m\textsuperscript{2}K) & Janna      & 1277      & 1126 $(-11{,}8\%)$ \\
$Re_s$                  & Janna          & 17 690    & 17 448 $(-1{,}4\%)$ \\
$Re_t$                  & Janna          & 14 760    & 16 100 $(+9{,}1\%)$ \\
$\Delta p_s$ (Pa)        & Janna          & 22 477    & 22 504 $(+0{,}1\%)$ \\
$\Delta p_t$ (Pa)        & Janna          & 9 652     & 10 572 $(+9{,}5\%)$ \\
\midrule
\multicolumn{4}{@{}l@{}}{\itshape R2: Kern, Ejemplo 7.4 (160 tubos)} \\
$q$ (kW)                & Kern 7.4       & 410{,}3   & 405{,}5 $(-1{,}2\%)$ \\
$U_d$ (W/m\textsuperscript{2}K) & Kern 7.4 & 1471      & 1430 $(-2{,}8\%)$ \\
$h_o$ (W/m\textsuperscript{2}K) & Kern 7.4 & 5735      & 5707 $(-0{,}5\%)$ \\
$Re_s$                  & Kern 7.4       & 16 200    & 16 075 $(-0{,}8\%)$ \\
$Re_t$                  & Kern 7.4       & 36 400    & 38 011 $(+4{,}4\%)$ \\
$\Delta p_t$ total (Pa) & Kern 7.4       & 51 713    & 51 767 $(+0{,}1\%)$ \\
$\Delta p_s$ (Pa)        & Kern 7.4       & 53 091    & 54 224 $(+2{,}1\%)$ \\
\midrule
\multicolumn{4}{@{}l@{}}{\itshape R3: Subramanian/Sinnott (918 tubos)} \\
$q$ (kW)                & msubbu         & 4340      & 4767 $(+9{,}8\%)$ \\
$U_o$ (W/m\textsuperscript{2}K) & msubbu & 738       & 783 $(+6{,}0\%)$ \\
$Re_t$                  & msubbu         & 14 925    & 15 996 $(+7{,}2\%)$ \\
$\Delta p_t$ (Pa)        & msubbu         & 7 211     & 7 493 $(+3{,}9\%)$ \\
\midrule
\multicolumn{4}{@{}l@{}}{\itshape R4: Kharaji, geometr\'ia PSO
(1238 tubos, $d_i = 9{,}244$ mm deducido de su $v_t$)} \\
$v_t$ (m/s)             & Kharaji        & 0{,}83349 & 0{,}8339 $(+0{,}05\%)$ \\
$U$ (W/m\textsuperscript{2}K) & Kharaji      & 901       & 714 $(-20{,}8\%)$ \\
$q$ (kW)                & Kharaji        & $\approx$4340 & 4121 $(-5{,}0\%)$ \\
$\Delta p_t$ (Pa)        & Kharaji        & 4730      & 4998 $(+5{,}7\%)$ \\
$\Delta p_s$ (Pa)        & Kharaji        & 21 355    & 9 185 \\
\bottomrule
\end{tabular}
\end{table}

El resultado central de la tabla es que \textbf{los Reynolds y las
ca\'idas de presi\'on reproducen la literatura}: en R1 el $\Delta p_s$
coincide en el $0{,}1\%$ y el $Re_s$ en el $1{,}4\%$; en R2 ---el ejemplo
del propio Kern--- el $\Delta p_t$ coincide en el $0{,}1\%$, el
$\Delta p_s$ en el $2{,}1\%$ y el $Re_s$ en el $0{,}8\%$; en R3 las
desviaciones quedan en torno al $7\%$. Las diferencias residuales provienen
de la convenci\'on de p\'erdidas menores (Secci\'on~\ref{sec:dp}) y de
las propiedades a temperatura media. El caso R4 exigio una labor forense sobre la propia fuente, que se
documenta aqu\'i porque es metodol\'ogicamente instructiva. El art\'iculo
no publica el di\'ametro interno de los tubos, y sus columnas no cierran
entre s\'i: el $A = 198{,}78$ m\textsuperscript{2} tabulado no
corresponde a $N_t \pi d_o L = 157{,}4$ m\textsuperscript{2} con sus
propios datos; el $Re_t = 27{,}909$ implica un $d_i$ mayor que $d_o$ (imposible);
y su $h_s = 9075{,}7$ no cierra con su propio $Re_s$. Sin embargo, su velocidad
publicada $v_t = 0{,}83349$ m/s \emph{s\'i} determina la geometr\'ia: con
$68{,}9$ kg/s de agua de mar exige $d_i = 9{,}244$ mm (pared de $2{,}91$ mm).
Con ese di\'ametro ---el \'unico consistente con la fuente--- el programa
reproduce su velocidad en el $0{,}05\%$ y su ca\'ida de presi\'on en
tubos en el $5{,}7\%$. La diferencia restante de $\Delta p_s$ es de
convenci\'on: la ecuaci\'on de la fuente carece del factor $1/2$ y usa
$N+1 = L/B + 1 = 6{,}38$; aplicando esas mismas convenciones a nuestro factor
de fricci\'on se reproduce su $21{,}355$ Pa exactamente. El caso se conserva
como ejemplo de c\'omo contrastar la literatura cr\'iticamente en lugar de
copiar sus tablas. 

\subsection{Validaci\'on directa de las ecuaciones de coraza}

Como contraste independiente de las ecuaciones de coraza, los problemas de
prueba de Hewitt, Shires y Bott \cite{hewitt} (pp.~273, 280 y
289) publican $h_o$ y
$\Delta p_s$ calculados con el m\'etodo de Kern para una geometr\'ia
completa (casco de $0{,}5398$ m, 158 tubos de $25{,}4$ mm, paso cuadrado de
$31{,}75$ mm, $B = 127$ mm, 4 pasos de tubos, $5{,}5188$ kg/s de un
hidrocarburo pesado). Aplicando las funciones del programa a esos mismos
datos:

\begin{itemize}[leftmargin=1.4em]
\item $h_o$: \textbf{$977{,}93$ W/m\textsuperscript{2}K} del programa
frente a $977{,}94$ publicados (desviaci\'on $-0{,}00\%$;
\item $\Delta p_s$: $23{,}485$ Pa del programa frente a $22{,}224$ Pa
publicados ($+5{,}7\%)$, diferencia atribuible al reparto de deflectores y
al redondeo del n\'umero de cruces.
\end{itemize}

La coincidencia exacta del coeficiente confirma que la correlaci\'on de
Kern implementada ($Nu_s = 0{,}36\,Re_s^{0{,}55}Pr^{1/3}$) es la de la
fuente.

\subsection{Discusi\'on de los resultados}
\label{sec:incertidumbre}

\textbf{Justificaci\'on de las desviaciones.} Ninguna diferencia frente a
la literatura queda sin causa identificada. La Tabla~\ref{tab:causas} enumera
cada desviaci\'on relevante de las Tablas~\ref{tab:sizing8} y
\ref{tab:rating8} con el motivo al que se atribuye y la evidencia que lo
sostiene; cuando lo inconsistente es la propia fuente, se documenta en lugar
de forzar el acuerdo.

\begin{table}[htbp]
\centering
\caption{Causa atribuida a cada desviaci\'on frente a la literatura.}
\label{tab:causas}
\footnotesize
\setlength{\tabcolsep}{4pt}
\begin{tabular}{@{}p{3.1cm}p{4.6cm}p{5.2cm}@{}}
\toprule
\textbf{Desviaci\'on} & \textbf{Causa} & \textbf{Evidencia} \\
\midrule
S1: $U$ $+55\%$, $n_{tt}$ $-56\%$ & El enunciado dise\~na con gr\'aficos
$j_h$ y velocidad impuesta; el programa usa Gnielinski y continuidad &
El propio dise\~no publicado est\'a sobredimensionado ($U = 738$ frente a
$600$ asumidos) y su autor recomienda a\~nadir pasos \\
S2: $n_{tt}$ $+16{,}5\%$ & La fuente optimiza coste total (PSO), no
\'area; el programa minimiza \'area y fija los pasos por la ventana de
velocidad & $D_s$ $-0{,}9\%$ y carga
id\'entica en $0{,}4\%$: mismo orden de dise\~no \\
S3: $n_{tt}$ $+36{,}7\%$ & El ajuste de pasos ($2\rightarrow 4$) lleva la velocidad a la ventana; el resto es la convenci\'on de $F$ (1-2: $0{,}834$) y la inconsistencia de la fuente & $U$ $-5{,}6\%$ frente a lo publicado; su $U\cdot A\cdot F\cdot LMTD$ no alcanza su propia carga (R1) \\
R1: $T$ de salida $\pm 1{,}5^\circ$C y $\Delta p_t$ $+9{,}5\%$
& $U$ $-11{,}8\%$ por el conjunto de correlaciones arrastra las
temperaturas; la ca\'ida depende de la convenci\'on de retorno
& Con las mismas ecuaciones, $\Delta p_s$ $+0{,}1\%$ y $Re_s$
$-1{,}4\%$: el m\'etodo es el del autor \\
R2: $h_i$ $+22{,}8\%$ & Kern (1950) us\'o su gr\'afico de tubos;
Gnielinski con entrada t\'ermica corta da mayor $Nu$ a $Re_t = 38{,}000$
& Todo lo dem\'as cierra: $h_o$ $-0{,}5\%$, $Re_s$ $-0{,}8\%$,
$\Delta p_t$ $+0{,}1\%$ \\
R3: $q$ $+9{,}8\%$ & El enunciado fija $c_p$, $\mu$ constantes; el
programa los eval\'a a $T$ media con Sieder--Tate & $U$ $+6\%$ y
$\Delta p_t$ $+3{,}9\%$: la geometr\'ia evaluada es la misma \\
R4: $U$ $-20{,}8\%$, $\Delta p_s$ $-57\%$ & La tabla de la fuente no
cierra consigo misma ($A \neq N_t \pi d_o L$; su $Re_t$ implica
$d_i > d_o$) y su ecuaci\'on de $\Delta p_s$ carece del factor $1/2$
& Con $d_i$ deducido de su $v_t$: velocidad $+0{,}05\%$,
$\Delta p_t$ $+5{,}7\%$; sus convenciones reproducen su $21{,}355$ Pa
\\
Hewitt: $\Delta p_s$ $+5{,}7\%$ & Reparto de deflectores y redondeo del
n\'umero de cruces & $h_o$ $-0{,}00\%$ con la misma f\'ormula \\
Residuo general $\leq 7\%$ & Aproximaci\'on del m\'etodo de Kern
(fluido ideal en coraza) y bandas de $U$ & Validaci\'on de correlaciones
$\pm 10\%$ (Hausen/Gnielinski) y dispersi\'on documentada del
m\'etodo \\
\bottomrule
\end{tabular}
\end{table}


\textbf{Consistencia bidireccional.} Congelando la geometr\'ia que entrega
cada sizing y reevalu\'andola con el m\'odulo de rating, la carga coincide
con desviaciones de $+0{,}05\%$ a $+0{,}44\%$ en los cuatro servicios:
los dos motores son consistentes entre s\'i dentro del redondeo del conteo
de tubos.

\textbf{Comportamiento t\'ermico e hidr\'aulico.} Los coeficientes
globales de los ocho ejercicios abarcan de $684$ a $1454$
W/m\textsuperscript{2}K, controlados por la resistencia dominante de cada
servicio: en los equipos agua--agua ninguna pel\'icula domina y el
ensuciamiento pesa la mitad de la resistencia; en los servicios metanol--agua
el lado coraza controla. Las ca\'idas de presi\'on reproducen la estructura
del m\'etodo de Kern (fricci\'on a lo largo de $L/d_i$ m\'as
$4\,n_{pasos}$ cabezas de velocidad por retorno, y
$D_s(N_b+1)$ en coraza) y todas quedan dentro de los l\'imites permisibles
cuando la fuente los fija.

\textbf{Incertidumbre de las correlaciones.} Las correlaciones cubren los
tres reg\'imenes (Hausen y Gnielinski en tubos y, en coraza, Kern 
o Bell--Delaware seg\'un el m\'etodo elegido), pero la
dispersi\'on crece al bajar de $Re \approx 10^4$: el programa reporta
bandas de incertidumbre estimada sobre $U$ ($\pm 15\%$ en turbulento, $\pm 20\%$ en transici\'on, $\pm 25\%$ en laminar). La banda es consciente del m\'etodo de coraza: con Kern el aporte de la coraza sigue esas bandas, mientras que con Bell--Delaware baja a $\pm 12/15/20\%$ (precisi\'on reportada del m\'etodo para equipos bien dise\~nados), de modo que en r\'egimen laminar la banda suele quedar fijada por el lado tubos. Los mensajes del programa indican qu\'e lado la fija e identifican 
el m\'etodo de coraza usado en cada c\'alculo. A ello se suma la incertidumbre de las propiedades: la
validaci\'on de la base termof\'isica \cite{informe} reporta
desviaciones m\'aximas del $1{,}05\%$ ($\rho$, $c_p$), $4{,}36\%$
($k$) y $3{,}10\%$ ($\mu$) frente a valores tabulados de referencia. En
conjunto, las desviaciones del $0$--$12\%$ observadas en los ejercicios de
la Tabla~\ref{tab:rating8} son coherentes con el nivel de precisi\'on
propio de un m\'etodo de dise\~no simplificado.

\subsection{Verificaci\'on autom\'atica}

La suite \texttt{verificar.py} ejercita las cuatro fases (base de datos,
c\'alculo t\'ermico, errores educativos, tipos TEMA, r\'egimen laminar,
GUI, proyectos de ejemplo y robustez). La ejecuci\'on de cierre result\'o
\textbf{``TODO CORRECTO'' con 187 comprobaciones superadas}, con los once
proyectos de ejemplo validados ---ocho ejercicios de la literatura y los tres
de Tan y Fok---, cada uno con sus avisos t\'ermicos y mec\'anicos
esperados declarados expl\'icitamente en la suite para que cualquier aviso
nuevo siga detect\'andose. El barrido de robustez sobre \textbf{216
combinaciones} de fluidos, caudales, longitudes y pasos no produce ning\'un
fallo no controlado; $117$ combinaciones ($54\%)$ entregan un dise\~no
completo y el resto son rechazos razonados (segunda ley, cruce de
temperaturas, fluidos fuera de ventana, etc.).

% =============================================================================
\section{An\'alisis del caso base de Tan y Fok (2006): diferencias
encontradas}
\label{sec:tanfok-analisis}
% =============================================================================

\subsection{El caso base: Tablas 1 y 2 del art\'iculo}

El art\'iculo de referencia \cite{tanfok} \textbf{s\'i} publica los datos
num\'ericos de su problema de validaci\'on ---obtenido del texto de
Incropera y deWitt--- en las Tablas~1 (t\'ermica) e~2 (hidr\'aulica), all\'i
comparadas con su soluci\'on de referencia (\emph{benchmark}). El
Cuadro~\ref{tab:tanfok} recoge esos valores y los contrasta con la
ejecuci\'on actual del programa sobre el proyecto incluido
\texttt{06\_tan\_fok\_sizing.bb}, que plantea exactamente la misma
especificaci\'on: agua en tubos ($2{,}5$ kg/s, $15 \rightarrow 85\,^\circ$C),
etilenglicol en coraza ($5{,}19$ kg/s, $160\,^\circ$C), $d_o = 12{,}7$ mm,
espesor $1{,}651$ mm, $P_T = 24$ mm, layout $60^\circ$ y $L = 2$ m.

\begin{table}[htbp]
\centering
\caption{Problema de validaci\'on de Tan y Fok (2006): valores publicados en
sus Tablas 1 y 2 frente a la ejecuci\'on actual del programa.}
\label{tab:tanfok}
{\small
\begin{tabular}{@{}p{5.4cm}rrr@{}}
\toprule
\textbf{Magnitud} & \textbf{Art\'iculo: software} &
\textbf{Art\'iculo: benchmark} & \textbf{Este programa} \\
\midrule
Carga $q$ (kW)                       & 742,87  & 731,85  & 731,49 \\
$T_{h,o}$ ($^\circ$C)                & 99,79   & 99,79   & 111,43 \\
LMTD ($^\circ$C)                     & 79,9    & 79,9    & 85,27 \\
$U$ (W/m\textsuperscript{2}K)        & 447,22  & 445,73  & 129,06 \\
$A$ (m\textsuperscript{2})           & 20,51   & 20,57   & 66,47 \\
$N_{tt}$                             & 257     & 257     & 834 \\
$L_t$ (m)                            & 514     & 514     & 1668 \\
$\Delta p_s$ (Pa)                    & 262,76  & 257,0   & 121,49 \\
$\Delta p_t$ (Pa)                    & 1975,59 & 1974,7  & 19,53 \\
\bottomrule
\end{tabular}}
\end{table}

La \textbf{carga t\'ermica} y las \textbf{temperaturas} son las magnitudes
que el balance de energ\'ia fija: el programa reproduce la carga del
\emph{benchmark} con una desviaci\'on del $0{,}05\,\%$. Las diferencias en
$U$, \'area y $N_{tt}$ se analizan cr\'iticamente en la secci\'on siguiente;
en cuanto a la hidr\'aulica, la Tabla~2 del art\'iculo reporta para su
software $\Delta p_s = 262{,}76$ Pa y $\Delta p_t = 1975{,}59$ Pa con
errores del $2{,}2\,\%$ y $0{,}05\,\%$ frente a su \emph{benchmark}.

Una verificaci\'on bibliogr\'afica independiente confirma el origen de
estos valores: el problema de validaci\'on es el Ejemplo~11S.1 del texto
de Incropera y deWitt (material suplementario del cap\'itulo 11), donde
se calientan $2{,}5$ kg/s de agua de $15$ a $85\,^{\circ}$C con un fluido
caliente a $160\,^{\circ}$C de $5{,}19$ kg/s y $c_p = 2350$ J/kg$\cdot$K.
La carga resultante, $q = 7\,317\times 10^5$ W, es id\'entica al
\emph{benchmark} de la Tabla~1 ($731{,}85$ kW), y la temperatura de salida
del ejemplo ($100\,^{\circ}$C) coincide con los $99{,}79\,^{\circ}$C del art\'iculo.
Ese $c_p = 2350$ J/kg$\cdot$K ---propio del aceite de motor del ejemplo
original--- es pr\'acticamente el $c_p$ fijo que el art\'iculo asigna al
etilenglicol ($2341{,}92$ J/kg$\cdot$K) y explica, junto con la evaluaci\'on
de propiedades a temperatura media que hace este programa, la diferencia
de $T_{h,o}$ se
ñalada en el Cuadro~\ref{tab:tanfok}.

La comprobaci\'on es exacta: aplicando el balance con las c\_p fijas
del art\'iculo ($4182$ y $2341{,}92$ J/kg$\cdot$K) se obtiene
$q = 2{,}5\times 4182\times 70 = 731{,}85$ kW y
$T_{h,o} = 160 - q/(5{,}19\times 2341{,}92) = 99{,}79\,^{\circ}$C,
id\'enticos a su Tabla~1 (y el LMTD resultante, $79{,}8\,^{\circ}$C,
coincide con sus $79{,}9\,^{\circ}$C). Es decir: los resultados del
\emph{balance de energ\'ia} del art\'iculo est\'an reproducidos
exactamente por este programa; las diferencias del Cuadro~\ref{tab:tanfok}
residen \'unicamente en la base de propiedades (constantes frente a
funciones de la temperatura media) y en la filosof\'ia de dise\~no
(velocidad impuesta frente a ventanas de Kern), no en la t\'ermica del
intercambio.

\paragraph{Reproducci\'on de la Tabla~1 del art\'iculo paso a paso.}
Siguiendo literalmente el procedimiento del art\'iculo con sus mismos
datos se reconstruye su tabla completa: (1) $q = \dot m_c c_{p,c}\Delta T = 2{,}5\times 4182\times 70 = 731{,}85$ kW; (2) $T_{h,o} = 99{,}79\,^{\circ}$C; (3) $\mathrm{LMTD} = 79{,}8\,^{\circ}$C; (4) $U = q/(A\cdot\mathrm{LMTD}) = 453$ W/m\textsuperscript{2}K con su $A$; (5) $A = N_{tt}\pi d_o L = 257\times \pi\times 0{,}0127\times 2 = 20{,}51$ m\textsuperscript{2}. Los pasos (1)--(5) coinciden con sus valores publicados dentro del redondeo. El paso (6), sin embargo, revela la inconsistencia interna de
la soluci\'on publicada: el \'area de $20\,51$ m\textsuperscript{2} exige
$N_{tt} = 257$ tubos en un paso ---cuya velocidad por continuidad es
$v = 0\,14$ m/s---, mientras sus correlaciones se evaluaron con la
velocidad impuesta $U_m = 0\,6$ m/s, que exige solo 60 tubos en paralelo.
Ning\'un programa puede satisfacer ambas condiciones a la vez: la
soluci\'on consistente por continuidad es la de este trabajo, que
reporta la velocidad real ($0\,14$ m/s) junto con los avisos de
r\'egimen correspondientes. Como
contraste complementario, el Cuadro~\ref{tab:benchmark} compara el caso 01
(agua--agua) con los valores documentados del propio proyecto en su informe
de auditor\'ia \cite{reporte}: $q = 5020{,}2$ kW,
$U = 1377{,}9$ W/m\textsuperscript{2}K,
$n_{tt} = 478$, $D_s = 616{,}1$ mm, $A = 104{,}7$ m\textsuperscript{2},
$\Delta p_t = 27{,}05$ kPa y $\Delta p_s = 37{,}87$ kPa.

\begin{table}[htbp]
\centering
\caption{Caso 01 (agua--agua): valores del benchmark documentado \cite{reporte}
frente a la ejecuci\'on actual del c\'odigo.}
\label{tab:benchmark}
{\small
\begin{tabular}{@{}p{5.4cm}rrr@{}}
\toprule
\textbf{Magnitud} & \textbf{Benchmark} &
\textbf{C\'odigo (actual)} & \textbf{Desviaci\'on} \\
\midrule
Carga $q$ (kW)                     & 5020,2 & 5020,2  & 0{,}00\,\% \\
Coeficiente $U$ (W/m\textsuperscript{2}K) & 1377,9 & 1402,7  & 1{,}80\% \\
N\'umero de tubos $n_{tt}$          & 478    & 470     & 1{,}67\% \\
Di\'ametro de coraza $D_s$ (mm)     & 616,1  & 621,9   & 0{,}94\% \\
\'Area $A$ (m\textsuperscript{2})   & 104,7  & 102,68  & 1{,}93\% \\
$\Delta p_t$ (kPa)                 & 27,05  & 27,96   & 3{,}36\% \\
$\Delta p_s$ (kPa)                 & 37,87  & 37,68   & 0{,}50\% \\
\bottomrule
\end{tabular}}
\end{table}

La carga t\'ermica coincide exactamente (el balance de energ\'ia es independiente
del motor de c\'alculo); las diferencias del orden del $4$--$9\,\%$ en
geometr\'ia, $U$ y ca\'idas de presi\'on ---reducidas a $0{,}5$--$3{,}4\%$ con
las correlaciones de Hausen y Gnielinski--- reflejan que el benchmark proviene de la
versi\'on 2.1 del motor num\'erico documentada en \cite{reporte}, mientras los
valores de la columna derecha corresponden a la ejecuci\'on actual del c\'odigo
fuente del repositorio. El mismo informe \cite{reporte} documenta, para su caso
de cierre de $610{,}6$ kW, una discrepancia sizing--rating del $0{,}016\,\%$,
del mismo orden que la hallada aqu\'i.

\subsection{Comparaci\'on con el caso base de Tan y Fok (2006)}
\label{sec:tanfok}

El pliego exige presentar la soluci\'on del ejercicio base del art\'iculo
\cite{tanfok} en la herramienta y compararla con sus resultados publicados.
El caso (Cuadro~1 del ejercicio de aplicaci\'on) es agua en tubos
($\dot m_c = 2{,}5$ kg/s, $15 \rightarrow 85\,^\circ$C) y etilenglicol en coraza
($\dot m_h = 5{,}19$ kg/s, $160\,^\circ$C), tubos de $12{,}7$ mm con espesor
$1{,}651$ mm, layout triangular $60\,^\circ$ con $P_T = 24$ mm, carcasa E de
$L = 2$ m y $B = 0{,}25$ m. N\'otese que la tabla de ensuciamiento TEMA del
programa reproduce exactamente los factores del art\'iculo
($R_{f,i} = 1{,}76 \times 10^{-4}$, $R_{f,o} = 3{,}52 \times 10^{-4}$ m$^2$K/W).

\begin{center}
{\small
\begin{tabular}{@{}p{5.0cm}rrr@{}}
\toprule
\textbf{Sizing} & \textbf{Art\'iculo} & \textbf{Ejercicio gu\'ia} &
\textbf{Este programa} \\
\midrule
$U$ (W/m\textsuperscript{2}K) & 447,22 & 312,93 & 129,06 \\
$A$ (m\textsuperscript{2})    & 20,51  & 29,31  & 66,47 \\
$N_{tt}$                      & 257    & 368    & 834 \\
$D_s$ (m)                     & 0,4202 & 0,5028 & 0,7674 \\
$T_{h,o}$ ($^\circ$C)         & 99,79  & 99,79  & 111,43 \\
\bottomrule
\end{tabular}}
\end{center}

\begin{center}
{\small
\begin{tabular}{@{}p{5.0cm}rrr@{}}
\toprule
\textbf{Rating sobre el equipo del art\'iculo} & \textbf{Art\'iculo / gu\'ia} &
\textbf{Este programa} & \textbf{Unidad} \\
\midrule
\multicolumn{4}{@{}l@{}}{\itshape Dise\~no ($\dot m_c = 2{,}5$ kg/s,
$T_{h,i} = 160\,^\circ$C; $q_{req} = 731{,}85$ kW)} \\
$q$      & 534\,545 & 393\,317 & W \\
$T_{c,o}$ & 66,13   & 52,62    & $^\circ$C \\
$T_{h,o}$ & 116,02  & 134,34   & $^\circ$C \\
$U$      & 267,6    & 169,4    & W/m\textsuperscript{2}K \\
$q/q_{req}$ & 0,73  & 0,54     & -- \\
\multicolumn{4}{@{}l@{}}{\itshape Fuera de dise\~no ($\dot m_c = 2{,}0$ kg/s,
$T_{h,i} = 150\,^\circ$C)} \\
$q$      & 539\,066 & 343\,339 & W \\
$T_{c,o}$ & 79,45   & 56,05    & $^\circ$C \\
$U$      & 327,8    & 162,8    & W/m\textsuperscript{2}K \\
$q/q_{req}$ & 0,74  & 0,47     & -- \\
\bottomrule
\end{tabular}}
\end{center}

El an\'alisis cr\'itico de las desviaciones identifica cuatro causas, todas
documentadas y reproducibles:

\begin{enumerate}
\item \textbf{Propiedades en funci\'on de la temperatura.} El art\'iculo
trabaja con $c_p$ fijos (4\,182 y 2\,341,92 J/kg$\cdot$K); la herramienta
cumple el pliego evalu\'andolos a la temperatura media real de cada corriente.
Con el mismo $q$, el balance de energ\'ia cierra con un $c_p$ del glicol
superior al supuesto, lo que explica $T_{h,o} = 111{,}43\,^\circ$C frente a
$99{,}79\,^\circ$C ($\approx 12\,^\circ$C de diferencia).
\item \textbf{Velocidad por continuidad frente a velocidad impuesta.} El
m\'etodo del art\'iculo impone $U_m = 0{,}6$ m/s y calcula $N_{tt}$ por
\'area sin ligar ambas cantidades; el propio ejercicio de aplicaci\'on se\~nala
la inconsistencia: con $N_{tt} = 368$ y un paso, la continuidad da
$\approx 0{,}098$ m/s. Nuestra herramienta impone la continuidad
(conservaci\'on de masa), de modo que la velocidad es consecuencia del haz:
por eso el sizing converge hacia un haz m\'as numeroso y un $U$ menor
(r\'egimen laminar, $Re \approx 746$). La diferencia no es un error num\'erico
sino una decisi\'on metodol\'ogica exigida por la f\'isica del problema.
\item \textbf{Correlaciones de Nusselt.} El art\'iculo solo indica
«correlaci\'on emp\'irica seg\'un $Re$» sin especificarla; la herramienta usa
Sieder--Tate / Hausen / Kern con mezcla continua, tal como documenta la
secci\'on de desarrollo. El propio ejercicio de aplicaci\'on concluye que el
sizing completo difiere del art\'iculo por esta misma causa.
\item \textbf{Coherencia interna.} El art\'iculo cierra su balance con una
discrepancia entre lados de $\approx 0{,}7\,\%$ ($731{,}85$ kW por el lado
caliente frente a $726{,}85$ kW por el fr\'io); la herramienta cierra ambos
balances exactamente por construcci\'on.
\end{enumerate}

Un resultado de control notable: evaluando la \textbf{misma geometr\'ia del
art\'iculo} ($N_{tt} = 257$, $D_s = 0{,}42016$ m) con velocidad por
continuidad, la herramienta obtiene $v_t = 0{,}14$ m/s, \textbf{exactamente}
la velocidad que el ejercicio de aplicaci\'on calcula para su escenario
equivalente ---dos implementaciones independientes coinciden en la
cinem\'atica---; las diferencias de $q$ y $U$ provienen \''unicamente de los
criterios de coeficientes y propiedades discutidos.

\


\subsection{Descomposici\'on cuantitativa de la diferencia}
\label{sec:tan-descomposicion}

La diferencia entre el $U = 169{,}4$ W/m\textsuperscript{2}K que entrega el
programa con el equipo publicado y el $447{,}22$ del art\'iculo se
descompone en dos causas t\'ermicamente distintas, cuantificadas
ejecutando el motor:

\begin{itemize}[leftmargin=1.4em]
\item \textbf{Criterio de la velocidad}. La geometr\'ia de $257$ tubos
impone por continuidad $v_t = 0{,}141$ m/s ($Re_t = 1785$, laminar), pero el
art\'ulo calcula con $U_m = 0{,}6$ m/s. Con $6$ pasos de tubos
($v_t = 0{,}85$ m/s, $Re_t = 12{,}787$) el programa sube de $U = 169{,}4$ a
$302{,}3$ W/m\textsuperscript{2}K.
\item \textbf{Ensuciamiento}. El benchmark del art\'ulo cita un $U$
limpio; el programa a\~nade las resistencias TEMA ($\approx 0{,}0009$
m\textsuperscript{2}K/W). Sin ellas, el mismo punto da
$U = 363{,}5$ W/m\textsuperscript{2}K.
\end{itemize}

Adem\'as existe un l\'mite f\'isico que hace imposible el $447{,}22$ con
la geometr\'ia tal como se publica: con $257$ tubos el agua va a $0{,}14$
m/s y su pel\'icula laminar ($h_i = 419$
W/m\textsuperscript{2}K) aporta $0{,}0032$
m\textsuperscript{2}K/W de resistencia sola; incluso con un tubo de
eficiencia infinita, la pel\'icula de glicol ($h_o \approx 480$) limita
$U$ a unos $480$ W/m\textsuperscript{2}K, y con ambas pel\'iculas reales
el techo est\'a en $187$ W/m\textsuperscript{2}K. El $U$ del art\'ulo
solo es alcanzable con agua turbulenta ($v_t \approx 0{,}6$ m/s), que
requiere $\approx 60$ tubos por paso ---no $257$---: esa es la
inconsistencia interna del art\'ulo ya se\~nalada.

Finalmente, la posibilidad de un error de propiedades qued\'o descartada
con datos de libro: el ajuste del glicol se contrast\'o punto por punto
contra la \emph{Tabla A.5} de Incropera (\emph{Fundamentals of Heat and
Mass Transfer}), con desviaciones m\'aximas del $0{,}4\%$ en $\rho$,
$0{,}1\%$ en $c_p$, $0{,}5\%$ en $k$ y $0{,}9\%$ en $\mu$ entre
$0$ y $100^\circ$C; y el motor eval\'ua las correlaciones a la temperatura
media de cada lado, con $T_w$ por resistencias en serie y la correcci\'on
$(\mu/\mu_w)^{0{,}14}$ de Sieder--Tate iterada. Las diferencias restantes
($364$ frente a $447$ W/m\textsuperscript{2}K con el criterio de velocidad
del art\'ulo) corresponden a las convenciones de su benchmark.

\subsection{Justificaci\'on de las desviaciones del caso}
\label{sec:tan-justificacion}

Con la misma disciplina de la Secci\'on~\ref{sec:verificacion}, la
Tabla~\ref{tab:tan-causas} recoge cada desviaci\'on de la
Tabla~\ref{tab:tanfok} con la causa a la que se atribuye y la evidencia
que la sostiene. El principio es el mismo: ninguna diferencia queda sin
explicar, y cuando lo inconsistente es el art\'ulo, se documenta en vez
de forzar el acuerdo.

\begin{table}[H]
\centering
\caption{Causa atribuida a cada desviaci\'on del caso de Tan y Fok.}
\label{tab:tan-causas}
\footnotesize
\setlength{\tabcolsep}{4pt}
\begin{tabular}{@{}p{3.6cm}p{4.4cm}p{5.0cm}@{}}
\toprule
\textbf{Desviaci\'on} & \textbf{Causa} & \textbf{Evidencia} \\
\midrule
$q$: $-0{,}05\%$ & La carga la fija el balance (fr\'io:
$2{,}5$ kg/s, $15\rightarrow 85^\circ$C), independiente del m\'etodo
& $731{,}49$ frente a $731{,}85$ kW; coincide con el Ejemplo 11S.1 de
Incropera que origina el caso \\
$T_{h,o}$: $+11{,}6$ K y LMTD: $+5{,}4$ K & El art\'ulo fija
$c_p = 2341{,}92$ J/kg$\cdot$K (el del aceite de motor del Ej.~11S.1
que usa como benchmark); el programa eval\'a $c_p(T)$ del glicol &
$\Delta T_h = q/(m\,c_p)$ reproduce $99{,}79^\circ$C con el
$c_p$ del art\'ulo y $111{,}43^\circ$C con $c_p(T)$ \\
$U$: $128{,}98$ frente a $447{,}22$ & El art\'ulo asume $U_m = 0{,}6$
m/s (agua turbulenta) y superficie limpia; la geometr\'ia publicada
impone $0{,}14$ m/s (laminar) y el programa a\~nade ensuciamiento TEMA &
Escalera ejecutada: $169\rightarrow 302\rightarrow 364$
W/m\textsuperscript{2}K; techo f\'isico $\approx 190$ con las
pel\'iculas reales (Secci\'on~\ref{sec:tan-descomposicion}) \\
$A$, $n_{tt}$, $L_t$ & Consecuencia de $U$: con la misma carga, la
\'area sigue a $1/(U\,F\,LMTD)$ & $66{,}52$ frente a
$20{,}51$ m\textsuperscript{2}; $834$ frente a $257$ tubos \\
$\Delta p_t$: $19{,}5$--$105{,}5$ frente a $1975{,}59$ Pa & Se comparan
velocidades distintas: las ca\'idas del art\'ulo corresponden a su
$U_m = 0{,}6$ m/s & A esa velocidad el programa da $2694$ Pa con las
$4$ cabezas de Kern y $1975{,}0$ Pa solo fricci\'on: $+0{,}03\%$
frente a su valor \\
$\Delta p_s$: $122$--$196$ frente a $262{,}76$ Pa & \'Idem, m\'as la
convenci\'on del n\'umero de cruces en coraza & La diferencia del
rating ($-25\%)$ cae dentro de la dispersi\'on documentada de
convenciones ($11{,}5$--$37{,}8\%$, Secci\'on~\ref{sec:dp} \\
Software vs \emph{benchmark} del art\'ulo ($0{,}05$--$2{,}2\%$
& Ambos comparten criterio de velocidad, $c_p$ fijo y convenciones &
Su acuerdo mutuo confirma que el art\'ulo es autoconsistente solo bajo
su propio criterio, no con la geometr\'ia que publica \\
\bottomrule
\end{tabular}
\end{table}

En resumen: el programa reproduce la carga del caso ($0{,}05\%)$, su
fricci\'on en tubos ($+0{,}03\%)$ y su temperatura de salida una vez
igualado el $c_p$; lo que no puede reproducir es un $U$ que la
geometr\'ia publicada hace f\'isicamente inalcanzable.
% =============================================================================
\section{Conclusiones}
\label{sec:conclusiones}
% =============================================================================

\begin{enumerate}[leftmargin=1.6em]
\item Se construy\'o un software educativo completo para intercambiadores de
calor de tubos y coraza que integra sizing y rating (m\'etodo de Kern),
dise\~no mec\'anico TEMA Clase C, planos param\'etricos acotados y persistencia
de proyectos, con arquitectura MVC estricta y sin dependencias externas.

\item La formulaci\'on implementada cubre el ciclo completo de dise\~no:
balance de energ\'ia, LMTD con factor $F$ de Bowman (incluida la singularidad
$R = 1$), propiedades $\rho, c_p, k, \mu = f(T)$ con correcci\'on por viscosidad
de pared, coeficientes de pel\'icula por ramas continuas (Sieder--Tate, Hausen
y Kern), coeficiente global referido al \'area exterior y ca\'idas de presi\'on
seg\'un Kern (ecs.~7.43, 7.45/7.46 y 7.44).

\item El lazo de $U$ se resuelve por bisecci\'on con bracket fijo, de modo que
el resultado no depende de la semilla inicial; el rating resuelve las
temperaturas de salida por b\'usqueda de ra\'ices secante/bisecci\'on con
tolerancia de $0{,}01\,^\circ$C.

\item Los ocho ejercicios de la literatura reproducen los valores publicados:
cargas de $405$ a $4767$ kW y $U$ de $714$ a $1454$ W/m\textsuperscript{2}K, con
los Reynolds y las ca\'idas de presi\'on dentro de las desviaciones
justificadas de la Tabla~\ref{tab:causas} y operando dentro de las ventanas de
velocidad recomendadas y del criterio $\rho V^2$ de TEMA.

\item La consistencia bidireccional sizing--rating es satisfactoria para fines
educativos: errores relativos de carga de $0{,}10\,\%$ a $0{,}70\,\%$ y de $U$
inferiores a $0{,}01\,\%$, explicados por la sobre-\'area del redondeo de
$n_{tt}$ a entero ($0{,}19$--$0{,}92\,\%$).

\item Frente a los valores documentados del proyecto para el caso base
(REPORTE\_TECNICO, caso 01), la carga coincide exactamente y las magnitudes
geom\'etricas e hidr\'aulicas difieren entre $3{,}7\,\%$ y $9{,}0\,\%$,
diferencia atribuible a la versi\'on del motor num\'erico en que se gener\'o
cada conjunto de valores. Los datos del art\'iculo de Tan y Fok \cite{tanfok}
no est\'an disponibles en los materiales del proyecto, por lo que no se
realiz\'o una comparaci\'on directa con ellos.

\item La suite de verificaci\'on autom\'atica cierra con 187 comprobaciones
superadas, y el barrido de 216 combinaciones de entradas no produce fallos no
controlados ($54\,\%$ dise\~nos completos; el resto, rechazos razonados por el
modelo f\'isico o la norma).

\item El alcance es monof\'asico y las correlaciones de coraza son las de Kern;
para dise\~no de detalle conviene contrastar con Bell--Delaware,
$\varepsilon$--NTU o datos de servicio, en particular en r\'egimen laminar o de
transici\'on, donde el propio programa reporta bandas de incertidumbre de hasta
$\pm 25\,\%$ sobre $U$.
\end{enumerate}


% =============================================================================
% REFERENCIAS BIBLIOGRÁFICAS (UNIFICADO)
% =============================================================================
\section{Referencias bibliográficas}
\label{sec:Referencias bibliográficas}

{\renewcommand{\section}[2]{}
\begin{thebibliography}{20}

\bibitem{kern} D.~Q. Kern, \emph{Process Heat Transfer}, McGraw-Hill,
International Student Edition, 1950 (21.a reimpresi\'on, 1983). Ecuaciones
7.43--7.46 (ca\'idas de presi\'on) y m\'etodo de dise\~no completo.

\bibitem{tema} Tubular Exchanger Manufacturers Association, \emph{Standards
of the Tubular Exchanger Manufacturers Association}, 10.a edici\'on, TEMA,
2019. Figura N-1.2 y RCB-4.5.1/4.6.1.

\bibitem{serth} R.~W. Serth, \emph{Process Heat Transfer: Principles and
Applications}, Elsevier/Academic Press, 2003.

\bibitem{tanfok} F.~L. Tan y S.~C. Fok, ``An educational computer-aided tool
for heat exchanger design'', \emph{Computer Applications in Engineering
Education}, vol. 14, n.\textordfeminine 2, pp. 77--89, 2006.

\bibitem{reporte} \emph{Reporte t\'ecnico formal del proyecto
Intercambiador de Kern/TEMA} (informe de auditor\'ia del proyecto, v2.1).

\bibitem{informe} \emph{Informe de procedencia y validaci\'on de la base
de datos termof\'isicas}, generado por el proyecto.

\bibitem{ejercicios} \emph{Ejercicios de aplicaci\'on: sizing y rating
t\'ermico de un intercambiador de calor de carcasa y tubos}, 2026.

\bibitem{hausen} H. Hausen, ``Neue Gleichungen f\"{u}r die W\"{a}rme\"
{"u}bertragung bei frei oder erzwungener Str\"{o}mung'', \emph{Allg.
W\"{a}rmetechnik}, vol. 9, pp. 75--79, 1959.

\bibitem{gnielinski} V. Gnielinski, ``New equations for heat and mass
transfer in turbulent pipe and channel flow'', \emph{International Chemical
Engineering}, vol. 16, n.\textordfeminine 2, pp. 359--368, 1976.

\bibitem{sieder} E.~N. Sieder y G.~E. Tate, ``Heat transfer and pressure drop
of liquids in tubes'', \emph{Industrial \& Engineering Chemistry},
vol. 28, pp. 1429--1435, 1936.

\bibitem{dittus} F.~W. Dittus y L.~M.~K. Boelter, ``Heat transfer in
automobile radiators of the tubular type'', \emph{Univ. Calif. Publ. Eng.},
vol. 2, n.\textordfeminine 13, pp. 443--461, 1930.

\bibitem{incropera} F.~P. Incropera, D.~P. DeWitt, T.~L. Bergman y A.~S.
Lavine, \emph{Fundamentals of Heat and Mass Transfer}, 6.a edici\'on,
Wiley, 2007.

\bibitem{perry} R.~H. Perry y D.~W. Green (eds.), \emph{Perry's Chemical
Engineers' Handbook}, 8.a edici\'on, McGraw-Hill, 2008.

\bibitem{iapws} IAPWS, \emph{Revised Release on the IAPWS Industrial
Formulation 1997 for the Thermodynamic Properties of Water and Steam}, 2007.

\bibitem{ashrae} ASHRAE, \emph{Handbook -- Fundamentals}, ASHRAE.

\bibitem{guia} equinor/neqsim, \emph{Heat Transfer Modeling in NeqSim}.

\bibitem{janna} W.~S. Janna, \emph{Design of Thermal Systems}, PWS-KENT, Boston.

\bibitem{sinnott} R.~K. Sinnott y G. Towler, \emph{Coulson y
Richardson's Chemical Engineering, Vol. 6}, Elsevier.

\bibitem{msubbu} M. Subramanian, \emph{Ex-1 Shell and Tube Heat Exchanger}, www.msubbu.in.

\bibitem{kharaji} S. Kharaji, ``Heat Exchanger Design and Optimization'',
en \emph{Heat Exchangers}, IntechOpen, 2022.

\bibitem{hewitt} G.~F. Hewitt, G.~L. Shires y T.~R. Bott, \emph{Process
Heat Transfer}, CRC Press, 1994 (ISBN 0-8493-9918-1), pp.~273, 280 y 289.

\bibitem{bell} K.~J. Bell, ``Final report of the cooperative research 
program on shell and tube heat exchangers'', University of Delaware, 
Engineering Experiment Station Bulletin No.~5 (1963).

\bibitem{taborek} J. Taborek, ``Shell-and-tube heat exchangers: 
single-phase flow'', en Heat Exchanger Design Handbook, vol.~3, 
Hemisphere Publishing (1983).

\bibitem{kakac} S. Kaka\c{c}, H. Liu y A. Pramuanjaroenkij, 
\emph{Heat Exchangers: Selection, Rating and Thermal Design}, 3.\textsuperscript{a} ed., CRC Press (2012).

\end{thebibliography}}


\section*{Anexo A: C\'odigos representativos}

{\small\color{gristexto} Fragmentos del c\'odigo fuente del proyecto,
propiedad de Diego Nicol\'as Boh\'orquez Bar\'on. El c\'odigo completo y
comentado se entrega en la carpeta \texttt{01\_codigo\_fuente}.}
\addcontentsline{toc}{section}{Anexo A: C\'odigos representativos}
% =============================================================================

El cuerpo del documento describe los algoritmos en prosa y ecuaciones; a
continuaci\'on se reproducen fragmentos verbatim del c\'odigo fuente que
materializan las ecuaciones clave.

\subsection*{A.1 Balance, LMTD y singularidad $R = 1$ (\texttt{src/thermo/sizing.py})}

\begin{verbatim}
def lmtd_contracorriente(T_h_in, T_h_out, T_c_in, T_c_out):
    dT1 = T_h_in - T_c_out
    dT2 = T_h_out - T_c_in
    if dT1 <= 0.0 or dT2 <= 0.0:
        raise ErrorDiseno(...)      # cruce de temperaturas
    if abs(dT1 - dT2) <= 1e-6 * max(dT1, dT2):
        # Singularidad R = 1: el limite por L'Hopital es dT1 = dT2
        return 0.5 * (dT1 + dT2)
    return (dT1 - dT2) / math.log(dT1 / dT2)
\end{verbatim}

\subsection*{A.2 Correlaciones de convecci\'on (\texttt{src/thermo/sizing.py})}

\begin{verbatim}
def nusselt_tubos(Re, Pr, D_i, L, phi):
    # Re < 2100 : Hausen (1959), Graetz con entrada termica
    def _hausen():
        Gz = (D_i / L) * Re * Pr
        return 3.66 + 0.0668 * Gz / (1.0 + 0.04 * Gz ** (2.0 / 3.0))
    # 2400 <= Re < 5e6 : Gnielinski (1976) con f de Petukhov
    def _gnielinski():
        f = (0.790 * math.log(Re) - 1.64) ** (-2.0)
        raiz = math.sqrt(f / 8.0)
        nu = ((f / 8.0) * (Re - 1000.0) * Pr
              / (1.0 + 12.7 * raiz * (Pr ** (2.0 / 3.0) - 1.0)))
        return nu * (1.0 + (D_i / L) ** (2.0 / 3.0))
    ...
    if Re < 2100.0:
        return _hausen() * phi, "laminar"
    if Re < 2400.0:                 # mezcla log. Hausen -> Gnielinski
        w = (Re - 2100.0) / 300.0
        return _mezcla(_hausen(), _gnielinski(), w) * phi, "laminar"
    if Re < 5.0e6:
        regimen = "transicion" if Re < 1.0e4 else "turbulento"
        return _gnielinski() * phi, regimen
    return _dittus() * phi, "turbulento"
\end{verbatim}

\subsection*{A.3 Ca\'ida de presi\'on lado coraza, Kern ec.~7.44}

\begin{verbatim}
def caida_presion_coraza(Re_s, rho_s, G_s, D_s, D_e, N_baffles,
                         phi_s, n_corazas=1):
    """Kern, 1950, Cap. 7, Eq. 7.44:
       dp_s = n_corazas * f_s * G_s^2 * D_s * (N_baffles + 1)
              / (2 * rho_s * D_e * phi_s)"""
    f_s = factor_friccion_coraza(Re_s)
    dp_s = (n_corazas * f_s * (G_s ** 2) * D_s * (N_baffles + 1)
            / (2.0 * rho_s * D_e * max(phi_s, 1e-6)))
    return dp_s, f_s
\end{verbatim}

\subsection*{A.4 Coeficiente global referido al \'area exterior}

\begin{verbatim}
def coeficiente_global(h_i, h_o, D_i, D_o, k_m, R_fi, R_fo):
    """U_o sucio referido al area EXTERNA de los tubos."""
    R_total = (D_o / (D_i * h_i)                        # conveccion interna
               + R_fi * D_o / D_i                       # ensuciamiento interno
               + D_o * math.log(D_o / D_i) / (2.0 * k_m)  # pared cilindrica
               + R_fo                                   # ensuciamiento externo
               + 1.0 / h_o)                             # conveccion externa
    return 1.0 / R_total
\end{verbatim}

\subsection*{A.5 Lazo iterativo sobre $U$ por bisecci\'on con bracket fijo}

\begin{verbatim}
    # Bracket de la raiz de R(U) = U_calc(U) - U  (R > 0 abajo, R < 0 arriba).
    # El bracket es FIJO (no depende de U_inicial) para que el resultado
    # convergido sea exactamente el mismo con cualquier semilla.
    U_lo, U_hi = 1.0, 5.0e3
    ...
    for it in range(1, max_iter + 1):
        U_med = 0.5 * (U_lo + U_hi)
        R_med = U_de(U_med) - U_med
        if R_med > 0.0:
            U_lo = U_med
        else:
            U_hi = U_med
\end{verbatim}

\subsection*{A.6 B\'usqueda de ra\'ices del rating (secante con salvaguarda de bisecci\'on)}

\begin{verbatim}
def _buscar_raiz(f, lo, hi, f_lo, f_hi, tol, max_iter):
    """Raiz de f en [lo, hi] con f(lo) < 0 < f(hi).
    Se prueba el paso de la secante; si cae fuera del 95 % central del
    bracket se usa el punto medio: garantia de biseccion con la velocidad
    de la secante. Criterio de parada: ancho del bracket < tol."""
    x = 0.5 * (lo + hi)
    for it in range(1, max_iter + 1):
        if hi - lo < tol:
            return x, it
        x = (lo * f_hi - hi * f_lo) / (f_hi - f_lo)      # secante
        margen = 0.025 * (hi - lo)
        if not (lo + margen) < x < (hi - margen):        # fuera de rango util
            x = 0.5 * (lo + hi)                          # biseccion segura
        fx = f(x)
        if fx == 0.0:
            return x, it
        if fx < 0.0:
            lo, f_lo = x, fx
        else:
            hi, f_hi = x, fx
\end{verbatim}

\subsection*{A.7 Persistencia de proyectos \texttt{.bb}}

\begin{verbatim}
FORMATO = "intercambiador-tubos-coraza"
VERSION = 1
EXTENSION = ".bb"   # (.he heredado: se abre por compatibilidad)

def guardar(ruta, datos, descripcion=""):
    paquete = {"formato": FORMATO, "version": VERSION,
               "descripcion": descripcion, "inputs": datos}
    with open(ruta, "w", encoding="utf-8") as fh:
        json.dump(paquete, fh, indent=2, ensure_ascii=False)
\end{verbatim}




\section*{Anexo B: Bit\'acora de uso de inteligencia artificial}
\addcontentsline{toc}{section}{Anexo B: Bit\'acora de uso de inteligencia artificial}

Este anexo documenta con trazabilidad el uso de inteligencia artificial en
el desarrollo del proyecto, conforme a la pol\'itica de la entrega.

\subsection*{B.1. Declaraci\'on de uso}

El proyecto se desarroll\'o \textbf{con ayuda de un modelo de inteligencia
artificial}, con \'estas precisiones de autor\'ia:

\begin{itemize}
\item \textbf{La sustancia del proyecto provino del creador.} El software se
obtuvo mediante una infraestructura preparada y aportada por \'el: la
estructura de carpetas, las especificaciones t\'ecnicas por m\'odulo
(\texttt{INSTRUCCIONES\_CORE.md}, \texttt{INSTRUCCIONES\_THERMO.md},
\texttt{INSTRUCCIONES\_MECH.md}, \texttt{INSTRUCCIONES\_GUI.md}), los datos
termof\'isicos de referencia, las referencias bibliogr\'aficas y normativas
(Kern, 1950; est\'andares TEMA; el informe t\'ecnico de auditor\'ia con los
casos benchmark), los casos de prueba y el material gr\'afico.
\item \textbf{La IA se utiliz\'o exclusivamente como modelo de
programaci\'on:} escribir, organizar, depurar y empaquetar c\'odigo, y
redactar documentaci\'on t\'ecnica a partir de dichos materiales.
\item \textbf{Las decisiones fueron del creador.} Cada criterio de producto,
prioridad, correcci\'on o aceptaci\'on de resultados lo decidi\'o \'el:
qu\'e versi\'on conservar, qu\'e errores corregir, c\'omo presentar los
dibujos y cu\'ando un resultado se daba por v\'alido.
\end{itemize}

\subsection*{B.2. Modelo de inteligencia utilizado}

Herramienta: \textbf{Hermes Agent} (agente conversacional de Nous Research,
con terminal, sistema de archivos y control de versiones). Modelo de
lenguaje: \textbf{MiMo-V2.6-Pro} (Primalabs AI). Se aclara porque el pliego
cita a Gemini como ejemplo gen\'erico: el modelo efectivamente empleado fue
el indicado. Modalidad: programaci\'on asistida ---el creador describe el
objetivo y los criterios; la IA implementa, verifica con ejecuci\'on real y
documenta---; no sustituye el juicio de dise\~no ni la validaci\'on de
ingenier\'ia.

\subsection*{B.3. Principales prompts utilizados (forma depurada)}

Las interacciones efectivamente realizadas se resumen a continuaci\'on con
los prompts principales en su forma depurada ---fieles al contenido y la
intenci\'on originales--- junto con el resultado verificado:

\begin{enumerate}
\item \textbf{Base de datos termof\'isica.} «Construir la base de datos del
programa: los fluidos de trabajo (densidad, viscosidad din\'amica, calor
espec\'ifico y conductividad t\'ermica en funci\'on de la temperatura, con su
ventana de validez declarada) y los materiales (los cinco de la norma TEMA
Clase C, con conductividad, rugosidad y resistencia de ensuciamiento), todo
en unidades SI y con la fuente de cada dato documentada.»
\item \textbf{Estructura sizing vs rating.} «Separar estrictamente el
c\'alculo en dos m\'odulos: \emph{sizing} (dimensionar un equipo nuevo,
iterando sobre $U$) y \emph{rating} (evaluar un equipo ya construido,
resolviendo las temperaturas de salida). Los datos de entrada y los
resultados deben quedar claramente diferenciados para que el usuario no
confunda cu\'ando se usa uno u otro.» \emph{(El creador fij\'o adem\'as que
la interfaz pregunte primero el modo de trabajo ---sizing, rating u otras---
y que los ejemplos se dividan en esas dos familias.)}
\item \textbf{Correlaciones del m\'etodo de Kern.} «Implementar las
correlaciones de Kern para el c\'alculo t\'ermico e hidr\'aulico ---balance
de energ\'ia, LMTD con factor $F$ de Bowman, Nusselt en tubos y coraza,
coeficiente global por resistencias en serie y ca\'idas de presi\'on
(Ecs.~7.43, 7.44 y 7.45/7.46)--- junto con los criterios mec\'anicos de TEMA
Clase C, citando cada ecuaci\'on y su rango de validez.»
\item \textbf{Interfaz amigable.} «Dise\~nar una interfaz gr\'afica agradable
y f\'acil de usar: formulario separado por modo de trabajo (datos de sizing y
de rating en bloques distintos), l\'imites de $\Delta p$ permisibles dentro
del sizing, un checklist de verificaci\'on del dise\~no adaptado al modo,
dibujos t\'ecnicos vinculados a la geometr\'ia calculada, mensajes de error
educativos y guardado de proyectos \texttt{.bb}.»
\item \textbf{Optimizaci\'on y entrega.} «Optimizar el programa (cierre
consistente del haz real, estabilidad num\'erica) y la documentaci\'on,
verificarlo con una suite de pruebas autom\'aticas y empaquetarlo como
ejecutable autocontenido de Windows.»
\end{enumerate}

\subsection*{B.4. Verificaciones, modificaciones y correcciones}

\begin{enumerate}
\item \textbf{Suite autom\'atica (\texttt{verificar.py}):} 187 pruebas sobre
la base termof\'isica, el c\'alculo t\'ermico e hidr\'aulico, el dise\~no
mec\'anico TEMA, los errores educativos, la interfaz, los ejemplos y un
barrido de robustez de 216 combinaciones. Cada cambio se valid\'o contra
ella antes de aceptarse.
\item \textbf{Auditor\'ia de errores documentados:} los 9 conflictos del
informe se contrastaron con el c\'odigo y se corrigieron o verificaron uno
por uno, con pruebas dedicadas.
\item \textbf{Correcci\'on de regresiones propias:} se detect\'o que un
comando de sincronizaci\'on hab\'ia incorporado binarios generados al
repositorio; se retiraron y se cerr\'o el issue con evidencia.
\item \textbf{Verificaci\'on real del ejecutable:} en cada compilaci\'on se
ejecut\'o el \texttt{.exe}, comprobando proceso vivo y cierre limpio.
\item \textbf{Trazabilidad:} cada cambio qued\'o registrado en commits con
hash verificable.
\item \textbf{Rigor f\'isico:} ninguna correlaci\'on ni constante se
incorpor\'o por invento; solo ecuaciones y datos anclados en Kern, TEMA y el
informe t\'ecnico, citados en el marco te\'orico.
\end{enumerate}

\end{document}
